% Билеты к первому коллоквиуму по курсу «Математическая логика» (ИТМО, осень 2026).
% Компилировать XeLaTeX / LuaLaTeX / tectonic (дважды — для навигации):
%   xelatex tickets.tex && xelatex tickets.tex
\documentclass[11pt,a4paper]{article}
\usepackage{fontspec}
\usepackage[english,russian]{babel}
% шрифт с кириллицей: первый найденный из списка
\IfFontExistsTF{Liberation Serif}{\setmainfont{Liberation Serif}}{%
\IfFontExistsTF{TeX Gyre Termes}{\setmainfont{TeX Gyre Termes}}{%
\IfFontExistsTF{FreeSerif}{\setmainfont{FreeSerif}}{\setmainfont{DejaVu Serif}}}}
\IfFontExistsTF{Liberation Mono}{\setmonofont{Liberation Mono}[Scale=MatchLowercase]}{\IfFontExistsTF{DejaVu Sans Mono}{\setmonofont{DejaVu Sans Mono}[Scale=MatchLowercase]}{}}
\usepackage[top=2.2cm,bottom=2.2cm,left=2cm,right=2cm,headheight=14pt]{geometry}
\usepackage{amsmath,amssymb,amsthm,stmaryrd,proof,enumitem}
\usepackage{tikz}
\usetikzlibrary{arrows.meta,positioning}
\usepackage[most]{tcolorbox}
\usepackage{fancyhdr}
\usepackage[unicode,colorlinks=true,linkcolor=blue!60!black,urlcolor=blue!60!black,bookmarksnumbered=true,bookmarksopen=true,pdftitle={Билеты к коллоквиуму 1 --- Математическая логика}]{hyperref}

\setlist{itemsep=1pt,topsep=3pt,parsep=0pt}
\setlength{\parindent}{0pt}
\setlength{\parskip}{4pt}

% ---------- обозначения ----------
\newcommand{\amp}{\mathbin{\&}}
\newcommand{\vd}{\vdash}
\newcommand{\nvd}{\nvdash}
\newcommand{\sem}[1]{\llbracket #1\rrbracket}
\newcommand{\True}{\text{И}}
\newcommand{\False}{\text{Л}}

\newcommand{\Gam}{\Gamma}
\newcommand{\cmpl}{\mathsf{c}}
\newcommand{\Vf}{\Vdash}
\newcommand{\cn}[1]{\llparenthesis #1\rrparenthesis}
\newcommand{\bbR}{\mathbb{R}}

% ---------- теоремные окружения ----------
\theoremstyle{definition}
\newtheorem*{dfn}{Определение}
\newtheorem*{thm}{Теорема}
\newtheorem*{lmm}{Лемма}
\newtheorem*{exm}{Пример}
\newtheorem*{rem}{Замечание}

\newtcolorbox{short}{colback=blue!5!white,colframe=blue!55!black,fonttitle=\bfseries,title={Коротко (что сказать в первую минуту)},boxrule=0.6pt,arc=2pt,left=4pt,right=4pt,top=2pt,bottom=2pt}
\newtcolorbox{warn}{colback=orange!7!white,colframe=orange!70!black,fonttitle=\bfseries,title={Типичные ловушки и дополнительные вопросы},boxrule=0.6pt,arc=2pt,left=4pt,right=4pt,top=2pt,bottom=2pt}

% ---------- навигация ----------
\newcounter{tk}
\newcommand{\NT}{31}% число билетов
\newcommand{\navbar}{%
  \ifnum\value{tk}>0
    \ifnum\value{tk}>1\hyperref[T:\the\numexpr\value{tk}-1\relax]{$\leftarrow$ предыдущий}\quad\fi
    \hyperref[toc]{оглавление}\quad\hyperref[map]{карта билетов}%
    \ifnum\value{tk}<\NT\quad\hyperref[T:\the\numexpr\value{tk}+1\relax]{следующий $\rightarrow$}\fi
  \fi}
\pagestyle{fancy}
\fancyhf{}
\fancyhead[L]{\small\ifnum\value{tk}>0 Билет \thetk\ из \NT\fi}
\fancyhead[R]{\small Математическая логика --- коллоквиум 1}
\fancyfoot[C]{\small\navbar}
\fancyfoot[R]{\small\thepage}
\renewcommand{\headrulewidth}{0.3pt}

% \ticket{номер вопроса в программе}{формулировка}
\newcommand{\ticket}[2]{%
  \clearpage
  \stepcounter{tk}%
  \phantomsection\label{T:\thetk}%
  \addcontentsline{toc}{subsection}{\protect\numberline{\thetk}#2}%
  \begin{center}
  \begin{tcolorbox}[colback=gray!8,colframe=gray!60,boxrule=0.5pt,arc=3pt,width=\linewidth]
  {\large\bfseries Билет \thetk}\hfill{\small вопрос программы: #1}\\[2pt]
  {\Large\bfseries #2}
  \end{tcolorbox}
  \end{center}
}

% \group{название} — титульная страница раздела
\newcommand{\group}[2]{%
  \clearpage
  \phantomsection
  \addcontentsline{toc}{section}{#1}%
  {\Large\bfseries #1}\par\medskip
  #2\par}

\begin{document}
\pagestyle{empty}
\begin{center}
{\LARGE\bfseries Математическая логика}\\[6pt]
{\Large Билеты к первому коллоквиуму}\\[4pt]
{\large ИТМО, группы M3232--M3239, осень 2026}\\[14pt]
\end{center}

Каждому вопросу из программы (\texttt{colloq1.tex}) соответствует один билет;
всего \NT\ билетов. В билете: формулировка вопроса, блок «Коротко»
(ядро ответа), полный ответ с определениями, теоремами и доказательствами
(или схемами доказательств, если на коллоквиуме требуется лишь формулировка),
блок типичных ловушек.

\textbf{Навигация.} В нижнем колонтитуле каждой страницы: переход к предыдущему/следующему
билету, к оглавлению и к «карте билетов». Все номера на карте, строки оглавления и
закладки PDF --- гиперссылки. Источники: лекции 1--4, конспекты 2011/2018/Чепелина,
теоретические домашние задания.

\bigskip
\phantomsection\label{toc}
\pagestyle{fancy}
\tableofcontents

\clearpage
\phantomsection\label{map}
\addcontentsline{toc}{section}{Карта билетов}
{\Large\bfseries Карта билетов}\par\medskip
\begin{center}
\renewcommand{\arraystretch}{1.5}
\begin{tabular}{|p{3.9cm}|p{11.6cm}|}
\hline
\textbf{1. Исчисление высказываний} & \hyperref[T:1]{1}\ парадоксы\quad \hyperref[T:2]{2}\ Гильберт\quad \hyperref[T:3]{3}\ метаязык\quad \hyperref[T:4]{4}\ язык ИВ\quad \hyperref[T:5]{5}\ оценка\quad \hyperref[T:6]{6}\ вывод\quad \hyperref[T:7]{7}\ корректность/полнота\quad \hyperref[T:8]{8}\ дедукция и полнота\\\hline
\textbf{2. Топология} & \hyperref[T:9]{9}\ определение\quad \hyperref[T:10]{10}\ примеры\quad \hyperref[T:11]{11}\ открытые, внутренность\quad \hyperref[T:12]{12}\ связность\quad \hyperref[T:13]{13}\ непрерывность\\\hline
\textbf{3. Виды вывода} & \hyperref[T:14]{14}\ гильбертов и натуральный вывод\\\hline
\textbf{4. Интуиционистское ИВ} & \hyperref[T:15]{15}\ чистое существование\quad \hyperref[T:16]{16}\ интуиционизм, BHK\quad \hyperref[T:17]{17}\ формальные логики\quad \hyperref[T:18]{18}\ решётки\quad \hyperref[T:19]{19}\ дистрибутивность\quad \hyperref[T:20]{20}\ подрешётка, M$_3$, N$_5$\quad \hyperref[T:21]{21}\ псевдодополнение\quad \hyperref[T:22]{22}\ алгебра Линденбаума\quad \hyperref[T:23]{23}\ полнота\quad \hyperref[T:24]{24}\ модели Крипке\quad \hyperref[T:25]{25}\ Крипке$\to$алгебры\quad \hyperref[T:26]{26}\ нетабличность\quad \hyperref[T:27]{27}\ разрешимость\\\hline
\textbf{5. Карри--Ховард} & \hyperref[T:28]{28}\ изоморфизм\\\hline
\textbf{6. Дизъюнктивность} & \hyperref[T:29]{29}\ гёделева алгебра\quad \hyperref[T:30]{30}\ гомоморфизм\quad \hyperref[T:31]{31}\ дизъюнктивность ИИВ\\\hline
\end{tabular}
\end{center}

\bigskip
\textbf{Сквозные договорённости.} Греческие буквы $\alpha,\beta,\gamma,\dots$ --- метапеременные
для формул; $\Gamma,\Delta$ --- списки формул; $\amp$ --- конъюнкция; $\sem{\cdot}$ --- оценка;
$\vd$ --- доказуемость (выводимость), $\vd_{\text{и}}$ и $\vd_{\text{к}}$ --- в интуиционистском
и классическом исчислении, если нужно различать; $\models$ --- общезначимость/следование;
$\True$, $\False$ --- «истина», «ложь».

% ======================================================================
\group{1. Исчисление высказываний}{Билеты 1--8: историческая мотивация, язык, семантика, доказательства, корректность и полнота, теорема о дедукции.\par
\hyperref[T:1]{1}\ \hyperref[T:2]{2}\ \hyperref[T:3]{3}\ \hyperref[T:4]{4}\ \hyperref[T:5]{5}\ \hyperref[T:6]{6}\ \hyperref[T:7]{7}\ \hyperref[T:8]{8}}

% ----------------------------------------------------------------------
\ticket{1(а)}{Парадокс брадобрея, парадокс Рассела (это разные парадоксы)}
\begin{short}
Брадобрей: «бреет всех, кто не бреется сам» --- противоречие показывает лишь, что такого брадобрея
\emph{не существует}. Рассел: $X=\{x\mid x\notin x\}$ --- противоречие в \emph{наивной теории множеств}, где любое свойство задаёт множество; пришлось
менять аксиоматику. Структура одна (диагональный аргумент), следствия разные.
\end{short}

\textbf{Контекст.} Ньютон и Лейбниц --- неформальное исчисление бесконечно малых (критика --- Беркли);
Коши (пределы), Вейерштрасс (вещественные числа), Кантор (теория множеств, 1875) формализовали анализ.
Парадокс Рассела (1901) подорвал основание: «Никто не изгонит нас из рая, который основал Кантор» (Гильберт).

\textbf{Брадобрей.} На острове живёт брадобрей, который бреет тех и только тех, кто не бреется сам.
Бреется ли он сам?
\begin{itemize}
\item Если бреет себя --- нарушено «бреет только тех, кто не бреется сам».
\item Если не бреет себя --- нарушено «бреет всех, кто не бреется сам».
\end{itemize}
Вывод: \emph{такого брадобрея нет} --- условие описывает несуществующий объект. Это не противоречие
в теории, а доказательство от противного: $\neg\exists b\,\forall x\,(\text{бреет}(b,x)\leftrightarrow\neg\text{бреет}(x,x))$.
Как формула это тавтология классической логики предикатов.

\textbf{Рассел.} В наивной теории множеств (свойство $\varphi(x)$ определяет множество
$\{x\mid\varphi(x)\}$) положим $X=\{x\mid x\notin x\}$. Тогда
\[
X\in X \iff X\notin X.
\]
Если $X\in X$, то $X$ обладает определяющим свойством, значит $X\notin X$. Если $X\notin X$, то $X$
обладает свойством и обязано принадлежать $X$. Противоречие выводится \emph{внутри теории}, т.\,е.
наивная теория противоречива.

\textbf{Различие.} В брадобрее мы вправе сказать «такого человека нет» --- понятие брадобрея
не входит в аксиомы. В теории множеств существование $X$ \emph{гарантировано} схемой неограниченного
выделения (понимания), так что вывод «$X$ не существует» противоречит самой аксиоме, а не только случаю.
Отсюда и сомнение: нет ли подобного противоречия в определении вещественных чисел, просто глубже скрытого?

\textbf{Выходы.} Аксиоматика Цермело--Френкеля: вместо $\{x\mid\varphi(x)\}$ допускается лишь
выделение $\{x\in A\mid\varphi(x)\}$ из уже существующего множества; аксиома фундирования запрещает $x\in x$;
теория типов Рассела (множество имеет тип выше своих элементов); интуиционизм Брауэра (отказ от «неконструктивных»
объектов --- множество $X$ нигде не предъявлено).

\begin{warn}
\begin{itemize}
\item Не говорите «брадобрей --- то же, что Рассел»: у брадобрея парадокса нет, есть теорема о несуществовании.
\item В программе Гильберта (билет~\hyperref[T:2]{2}) ключевая тревога --- противоречивость самой формализации математики.
\item Аналогия с программированием: класс \texttt{IntList} с полем \texttt{next} ссылается сам на себя --- «множество, принадлежащее себе» само по себе не страшно; страшна неограниченная схема комбинаторного построения.
\end{itemize}
\end{warn}

% ----------------------------------------------------------------------
\ticket{1(б)}{Программа Гильберта}
\begin{short}
Полностью формализовать математику и доказать финитными средствами: полноту, непротиворечивость,
консервативность и разрешимость формализации. Теоремы Гёделя о неполноте (1930--31) показали, что в полном
объёме программа невыполнима.
\end{short}

\textbf{Формулировка.} Гильберт (доклады в Гамбурге, 1921): <<Neubegründung der Mathematik>>.
Программа состоит из пунктов:
\begin{enumerate}
\item \textbf{формализация} всей математики (язык, аксиомы, правила вывода);
\item \textbf{полнота} формализации: любой истинный факт доказуем;
\item \textbf{непротиворечивость}: из аксиом нельзя вывести противоречие, причём доказать это
  «надёжными» (финитными) средствами;
\item \textbf{консервативность}: любое доказательство утверждения о «реальных» (финитных) объектах, использующее
  «идеальные» объекты (актуальная бесконечность и т.\,п.), можно переписать без них;
\item \textbf{разрешимость}: есть алгоритм, проверяющий истинность любого математического утверждения
  (Entscheidungsproblem).
\end{enumerate}

\textbf{Судьба программы.}
Теоремы Гёделя о неполноте формальной арифметики (1930/31) показали: в непротиворечивой теории, содержащей
арифметику, есть истинные недоказуемые утверждения (первая теорема), а непротиворечивость такой теории
не доказуема её собственными средствами (вторая). Неразрешимость исчисления предикатов
(Чёрч, Тьюринг) закрыла пункт о разрешимости. Тем не менее методы программы породили теорию доказательств,
теорию моделей, теорию вычислимости.

\textbf{В курсе.} Мы следуем программе «в некоторой степени»: формализуем фрагменты (исчисление высказываний,
предикатов, формальная арифметика, теория множеств) и изучаем проблемы, возникающие при формализации:
корректность, полнота, непротиворечивость, разрешимость.

\begin{warn}
\begin{itemize}
\item Различайте полноту \emph{теории} (синтаксическую/семантическую) и полноту \emph{исчисления} (теорема Гёделя о полноте --- про общезначимость; Гёдель о неполноте --- про арифметику).
\item Исчисление высказываний программе удовлетворяет полностью: оно корректно, полно, непротиворечиво и разрешимо (таблицы истинности).
\end{itemize}
\end{warn}

% ----------------------------------------------------------------------
\ticket{1(в)}{Предметный язык и язык исследователя (метаязык). Соглашения об обозначениях. Схемы формул}
\begin{short}
Предметный язык --- формальный язык, тексты на котором изучаем; метаязык --- язык, на котором
изучаем (русский, математика, грамматики). Метапеременные $\alpha,\beta,\dots$ обозначают формулы; схема
формул --- «формула с метапеременными», задаёт бесконечное семейство формул подстановкой.
\end{short}

\textbf{Два языка.}
\begin{itemize}
\item \emph{Предметный язык} --- формальный язык исчисления; его выражения (формулы, доказательства) --- объекты изучения.
\item \emph{Метаязык} (язык исследователя) --- язык, на котором мы формулируем утверждения о предметном языке: «формула $\alpha$ доказуема», «если $\Gamma\vd\alpha$, то\dots».
\end{itemize}
Пример из программирования: предметный язык --- Haskell, метаязык --- русский, язык грамматик, блок-схем.
Исчисление определяется предметным языком, \emph{теорией моделей} (как оценивать формулы) и
\emph{теорией доказательств} (аксиомы, правила вывода).

\textbf{Смешение уровней опасно.} Символ $\vd$ и $\models$ принадлежат метаязыку, $\to$ --- предметному:
$\alpha\to\beta$ --- формула, $\alpha\vd\beta$ --- утверждение о выводимости. (Теорема о дедукции связывает их.)

\textbf{Соглашения об обозначениях (метаязык).}
\begin{itemize}
\item Метапеременные $\alpha,\beta,\gamma,\dots$ --- любые формулы; одинаковые буквы означают одинаковые формулы, разные --- не обязательно разные.
\item Метапеременные $X,Y_n,Z'$ --- пропозициональные переменные; $A,B,C_{1234}$ --- сами пропозициональные переменные.
\item Большие греческие $\Gamma,\Delta,\Sigma,\Pi$ --- списки формул (возможно, пустые): $\Gamma,\Delta,\zeta\vd\alpha$.
\item Упрощение записи: приоритеты $\neg,\ \amp,\ \vee,\ \to$ (по убыванию), конъюнкция и дизъюнкция левоассоциативны,
  импликация правоассоциативна:
  \[
  (A\to B)\amp Q\vee(\neg B\to B\to C)\vee(C\to C\to A)
  \]
  означает $((((A\to B)\amp Q)\vee((\neg B\to B)\to C))\vee(C\to(C\to A)))$.
  При подстановке выражений вместо метапеременных они берутся в скобки.
\end{itemize}

\textbf{Схемы формул.}
\begin{dfn}
\emph{Схема высказываний} --- строка, строящаяся по правилам для формул, но вместо
пропозициональных переменных допускаются метапеременные (греческие буквы).
Формула $\sigma$ \emph{строится по схеме} $\textit{Ш}$, если существует замена метапеременных
$\textit{ч}_1,\dots,\textit{ч}_n$ схемы на формулы $\varphi_1,\dots,\varphi_n$ с
$\sigma=\textit{Ш}[\textit{ч}_1:=\varphi_1]\dots[\textit{ч}_n:=\varphi_n]$.
\end{dfn}
\begin{exm}
Схема $A\to\alpha\vee B\vee\alpha$ задаёт $A\to X\vee B\vee X$ и $A\to(M\to N)\vee B\vee(M\to N)$, но
\emph{не} задаёт $A\to X\vee B\vee Y$ (все вхождения $\alpha$ заменяются одинаково) и не задаёт
формулу, где нарушена структура скобок.
\end{exm}
Аксиомы исчисления задаются \emph{схемами аксиом}: схема $\alpha\amp\beta\to\alpha$ порождает и $A\amp B\to A$, и $B\amp A\to B$.
Иначе пришлось бы выписывать бесконечно много отдельных аксиом.

\begin{warn}
\begin{itemize}
\item Схема не формула: нельзя спрашивать «доказуема ли схема» --- спрашивают, \emph{все ли} формулы по схеме доказуемы/общезначимы (общезначимость схемы аксиом).
\item Разные метапеременные могут принять одно значение: $\alpha\to\beta\to\alpha$ при $\alpha=\beta=A$ даёт $A\to A\to A$.
\end{itemize}
\end{warn}

% ----------------------------------------------------------------------
\ticket{1(г)}{Язык исчисления высказываний}
\begin{short}
Формула --- пропозициональная переменная либо $\neg\alpha$, $\alpha\amp\beta$, $\alpha\vee\beta$, $\alpha\to\beta$.
Приоритеты: $\neg>\amp>\vee>\to$; $\amp,\vee$ левоассоциативны, $\to$ правоассоциативна.
\end{short}

\textbf{Грамматика.} Основа предметного языка классического исчисления высказываний --- \emph{высказывание} (формула):
\begin{align*}
\langle\text{формула}\rangle ::=\ & \langle\text{переменная}\rangle\ \mid\ (\neg\langle\text{формула}\rangle)\ \mid\
(\langle\text{ф.}\rangle\amp\langle\text{ф.}\rangle)\ \mid\ (\langle\text{ф.}\rangle\vee\langle\text{ф.}\rangle)\ \mid\
(\langle\text{ф.}\rangle\to\langle\text{ф.}\rangle)\\
\langle\text{переменная}\rangle ::=\ & A\mid B\mid\dots\mid Z\ \text{ (с индексами и штрихами: } A,\ B',\ C_{1234})
\end{align*}
Атомарное высказывание --- пропозициональная переменная; остальные --- составные. Для конъюнкции допускают
запись $\wedge$, для импликации $\supset$.

\textbf{Сокращения.} Приоритет по убыванию: отрицание, конъюнкция, дизъюнкция, импликация.
Ассоциативность: левая для $\amp$ и $\vee$, правая для $\to$. Тем самым
\[
A\to B\to C \equiv A\to(B\to C),\qquad A\vee B\vee C\equiv(A\vee B)\vee C,\qquad
\neg A\amp B\vee C\equiv((\neg A)\amp B)\vee C.
\]
Пример: $(((A\to B)\vee(B\to C))\vee(C\to A))$.

\textbf{Выбор связок.} Четыре связки выбраны для удобства; в других курсах берут иной базис
(например $\to,\bot$), и язык при необходимости расширяют ($\bot$ вместо $\neg$ в натуральном выводе, билет~\hyperref[T:14]{14}).

\textbf{Дерево разбора.} Каждая формула однозначно разбирается в дерево с листьями-переменными и
внутренними вершинами-связками; подформула --- поддерево. На этом строятся индукция по структуре формулы и рекурсивное определение оценки.

\begin{warn}
\begin{itemize}
\item $A\to B\to C$ \emph{не} равно $(A\to B)\to C$ --- это разные формулы: при $A=B=C=\False$ первая истинна ($\False\to(\False\to\False)=\True$), вторая ложна ($(\False\to\False)\to\False=\False$).
\item Отрицание --- одноместная, остальные --- двуместные.
\end{itemize}
\end{warn}

% ----------------------------------------------------------------------
\ticket{1(д)}{Оценка высказываний, общезначимость, следование}
\begin{short}
Модель $f:\mathcal{P}\to\{\True,\False\}$ рекурсивно задаёт $\sem{\alpha}$. $\models\alpha$ --- истинна при всех оценках;
$\Gamma\models\alpha$ --- истинна при всех оценках, делающих истинными все $\gamma_i$.
\end{short}

\textbf{Теория моделей.} Фиксируем $V=\{\True,\False\}$ (истина, ложь). \emph{Оценка (интерпретация, модель)}
пропозициональных переменных --- функция $f:\mathcal{P}\to V$. Записываем
$\sem{\alpha}^{X_1:=v_1,\dots,X_n:=v_n}$ --- значение формулы при указанных значениях (это метаязык, полагаемся на здравый смысл).
Рекурсивно:
\begin{itemize}
\item $\sem{X}=f(X)$,\quad $\sem{X}^{X:=a}=a$;
\item $\sem{\neg\alpha}=\False$, если $\sem{\alpha}=\True$, иначе $\True$;
\item $\sem{\alpha\amp\beta}=\True$, если $\sem{\alpha}=\sem{\beta}=\True$, иначе $\False$;
\item $\sem{\alpha\vee\beta}=\False$, если $\sem{\alpha}=\sem{\beta}=\False$, иначе $\True$;
\item $\sem{\alpha\to\beta}=\False$, если $\sem{\alpha}=\True$ и $\sem{\beta}=\False$, иначе $\True$.
\end{itemize}

\begin{dfn}
\begin{itemize}
\item $\alpha$ \emph{общезначима} (тавтология), $\models\alpha$, если истинна при любой оценке.
\item $\gamma_1,\dots,\gamma_n\models\alpha$ (\emph{следствие}), если $\alpha$ истинна при любой оценке, при которой истинны $\gamma_1,\dots,\gamma_n$.
\item \emph{Выполнима} --- истинна при какой-нибудь оценке; \emph{невыполнима} --- не истинна ни при какой; \emph{опровержима} --- ложна при какой-нибудь оценке.
\end{itemize}
\end{dfn}

\textbf{Примеры.} $A\to A$ --- тавтология: $\sem{A\to A}^{A:=\True}=\sem{A\to A}^{A:=\False}=\True$.
$A\to\neg A$ тавтологией не является: $\sem{A\to\neg A}^{A:=\True}=\False$ (но выполнима: при $A=\False$ истинна).

\textbf{Таблицы.}
\[
\begin{array}{cc|cccc}
\alpha&\beta&\neg\alpha&\alpha\amp\beta&\alpha\vee\beta&\alpha\to\beta\\\hline
\False&\False&\True&\False&\False&\True\\
\False&\True&\True&\False&\True&\True\\
\True&\False&\False&\False&\True&\False\\
\True&\True&\False&\True&\True&\True
\end{array}
\]
Проверка общезначимости любой формулы с $n$ переменными --- перебор $2^n$ строк, поэтому
общезначимость \emph{разрешима} (хотя и экспоненциально).

\textbf{Классическая и другие оценки.} Описанная оценка называется классической. Позднее мы
берём другие множества значений (открытые множества пространства, элементы решётки, миры Крипке),
но классическая подразумевается, если не оговорено иное.

\begin{warn}
\begin{itemize}
\item $\Gamma\models\alpha$ \emph{не} то же, что $\models\Gamma\to\alpha$ для бесконечного $\Gamma$; для конечного $\Gamma=\{\gamma_1,\dots,\gamma_n\}$ эквивалентно $\models\gamma_1\to\dots\to\gamma_n\to\alpha$.
\item Знак $\models$ --- метаязык; формула $\alpha\to\beta$ --- предметный язык.
\end{itemize}
\end{warn}

% ----------------------------------------------------------------------
\ticket{1(е)}{Доказуемость, гипотезы (контекст), выводимость}
\begin{short}
Вывод --- конечная последовательность формул, каждая из которых аксиома, гипотеза или получена по Modus Ponens
из предыдущих. $\vd\alpha$ --- есть вывод из пустого контекста; $\Gamma\vd\alpha$ --- вывод с гипотезами $\Gamma$.
\end{short}

\textbf{Аксиомы.} Десять схем:
\begin{align*}
&(1)\ \alpha\to\beta\to\alpha &&(6)\ \alpha\to\alpha\vee\beta\\
&(2)\ (\alpha\to\beta)\to(\alpha\to\beta\to\gamma)\to(\alpha\to\gamma) &&(7)\ \beta\to\alpha\vee\beta\\
&(3)\ \alpha\to\beta\to\alpha\amp\beta &&(8)\ (\alpha\to\gamma)\to(\beta\to\gamma)\to(\alpha\vee\beta\to\gamma)\\
&(4)\ \alpha\amp\beta\to\alpha &&(9)\ (\alpha\to\beta)\to(\alpha\to\neg\beta)\to\neg\alpha\\
&(5)\ \alpha\amp\beta\to\beta &&(10)\ \neg\neg\alpha\to\alpha
\end{align*}
Все формулы, получаемые подстановкой в эти схемы, --- аксиомы. Единственное правило вывода --- \emph{Modus Ponens}
(Теофраст, IV--III вв. до н.э.):
\[
\infer{\beta}{\alpha & \alpha\to\beta}
\]

\begin{dfn}[доказательство]
\emph{Доказательство} (вывод) --- конечная последовательность формул $\delta_1,\dots,\delta_n$, где каждая $\delta_i$ либо
аксиома, либо получена по M.P. из $\delta_j$ и $\delta_k\equiv\delta_j\to\delta_i$ при $j,k<i$.
Доказательство формулы $\alpha$ --- такое, что $\delta_n\equiv\alpha$; тогда $\alpha$ \emph{доказуема}, $\vd\alpha$.
\end{dfn}
\begin{dfn}[вывод из гипотез]
Вывод $\alpha$ из гипотез $\gamma_1,\dots,\gamma_k$ --- последовательность $\delta_1,\dots,\delta_n$, где каждая $\delta_i$ либо аксиома,
либо M.P. из предыдущих, либо одна из гипотез. $\alpha$ \emph{выводима из} $\gamma_1,\dots,\gamma_k$: $\gamma_1,\dots,\gamma_k\vd\alpha$.
Гипотезы называют также допущениями, предположениями, контекстом.
\end{dfn}
Доказательство --- вывод из пустого контекста.

\textbf{Пример (доказательство $\vd A\to A$).}
\begin{align*}
&(1)\ A\to(A\to A) &&\text{сх.\ акс.\ 1}\\
&(2)\ (A\to(A\to A))\to(A\to((A\to A)\to A))\to(A\to A) &&\text{сх.\ акс.\ 2}\\
&(3)\ (A\to((A\to A)\to A))\to(A\to A) &&\text{M.P.\ 1,2}\\
&(4)\ A\to((A\to A)\to A) &&\text{сх.\ акс.\ 1}\\
&(5)\ A\to A &&\text{M.P.\ 4,3}
\end{align*}

\textbf{Свойства.} (а) Слабление: из $\Gamma\vd\alpha$ следует $\Gamma,\Delta\vd\alpha$ (и можно переставлять гипотезы).
(б) Если $\Gamma\vd\gamma$, $\Delta\vd\delta$ и $\gamma,\delta\vd\alpha$, то $\Gamma,\Delta\vd\alpha$ (склеиваем выводы). (в) Гипотеза выводима: $\Gamma,\alpha\vd\alpha$.

\textbf{Схемы и бесконечность.} Схемы аксиом и правила вывода --- удобство задания; считаем, что аксиом и правил
бесконечно много (подстановка всех формул).

\begin{warn}
\begin{itemize}
\item $\gamma_1,\dots,\gamma_k\vd\alpha$ --- про \emph{выводимость}, $\models$ --- про истинность. Их связь --- теоремы о корректности и полноте (билет~\hyperref[T:7]{7}).
\item Вывод конечен; поэтому в выводе используются лишь конечно много гипотез.
\end{itemize}
\end{warn}

% ----------------------------------------------------------------------
\ticket{1(ж)}{Корректность, полнота, противоречивость и непротиворечивость (эквивалентные формулировки)}
\begin{short}
Корректность: $\vd\alpha\Rightarrow\models\alpha$. Полнота: $\models\alpha\Rightarrow\vd\alpha$. Теория противоречива, если доказуемо $A\amp\neg A$;
эквивалентно: доказуема любая формула; доказуемы и $\alpha$, и $\neg\alpha$ для какого-то $\alpha$.
\end{short}

\begin{dfn}
Теория \emph{корректна}, если любое доказуемое утверждение общезначимо: $\vd\alpha$ влечёт $\models\alpha$.
Теория (семантически) \emph{полна}, если любое общезначимое утверждение доказуемо: $\models\alpha$ влечёт $\vd\alpha$.
\end{dfn}

\begin{thm}[корректность КИВ]
Если $\vd\alpha$, то $\models\alpha$.
\end{thm}
\begin{proof}
Индукция по длине вывода $\delta_1,\dots,\delta_n$ (доказываем общезначимость каждой $\delta_i$).
\begin{itemize}
\item \emph{Аксиома}: все схемы аксиом общезначимы (проверка таблицами; например, для схемы 9 переберём $\sem{\alpha},\sem{\beta}$: значение $(\alpha\to\beta)\to(\alpha\to\neg\beta)\to\neg\alpha$ всегда $\True$).
\item \emph{Modus Ponens}: пусть $\delta_j$ и $\delta_j\to\delta_i$ общезначимы. Фиксируем оценку: $\sem{\delta_j}=\sem{\delta_j\to\delta_i}=\True$. По таблице импликации $\sem{\delta_i}=\False$ возможно лишь если $\sem{\delta_j}=\False$ или $\sem{\delta_j\to\delta_i}=\False$ --- невозможно. Значит $\sem{\delta_i}=\True$.\qedhere
\end{itemize}
\end{proof}

\begin{thm}[полнота КИВ]
Если $\models\alpha$, то $\vd\alpha$. (Доказательство --- билет~\hyperref[T:8]{8}.)
\end{thm}
Вместе: $\vd\alpha\iff\models\alpha$. Следствие: если неверно $\vd\alpha$, то есть контрпример (оценка, опровергающая $\alpha$). Кроме того, из корректности и полноты следует \emph{разрешимость}:
доказуемость проверяется таблицей истинности.

\begin{dfn}[противоречивость]
Теория \emph{противоречива}, если доказуемо $A\amp\neg A$ (для некоторой формулы $A$).
\end{dfn}
\begin{thm}[эквивалентные формулировки; верно для КИВ и ИИВ]
Следующие условия эквивалентны:
\begin{enumerate}
\item теория противоречива ($\vd A\amp\neg A$ для некоторой $A$);
\item найдётся $\alpha$ такая, что $\vd\alpha$ и $\vd\neg\alpha$;
\item любая формула доказуема;
\item доказуема ложь (в ИИВ: $\bot$ или $A\amp\neg A$ при любом $A$).
\end{enumerate}
Следовательно, теория \emph{непротиворечива} $\iff$ существует недоказуемая формула.
\end{thm}
\begin{proof}[Схема]
(1)$\Rightarrow$(2): из $A\amp\neg A$ по аксиомам 4, 5 получаем $A$ и $\neg A$.
(2)$\Rightarrow$(3): формула $\alpha\to\neg\alpha\to\beta$ доказуема (принцип взрыва, в КИВ --- из аксиомы 10; в ИИВ --- аксиома 10и); два раза M.P.
(3)$\Rightarrow$(1): в частности доказуема $A\amp\neg A$.
\end{proof}
\textbf{Непротиворечивость КИВ} следует из корректности: если бы было $\vd A\amp\neg A$, то $A\amp\neg A$ была бы тавтологией, но она всегда ложна. Непротиворечивость ИИВ --- из непротиворечивости КИВ (ИИВ слабее).

\begin{warn}
\begin{itemize}
\item Определения полноты бывают разные (синтаксическая: $\vd\alpha$ или $\vd\neg\alpha$ для замкнутых $\alpha$; семантическая). В КИВ имеем семантическую полноту; синтаксическая полнота к КИВ неприменима (переменные не замкнуты).
\item В теории предикатов/арифметики корректность не гарантирует полноты (теоремы Гёделя о неполноте).
\end{itemize}
\end{warn}

% ----------------------------------------------------------------------
\ticket{1(з)}{Теорема о дедукции (формулировка). Теорема о полноте (формулировка)}
\begin{short}
Теорема о дедукции: $\Gamma,\alpha\vd\beta\iff\Gamma\vd\alpha\to\beta$. Теорема о полноте: $\models\alpha\Rightarrow\vd\alpha$.
Дополнительно приводим идею доказательств.
\end{short}

\begin{thm}[о дедукции, Эрбран, 1930]
При любых $\Gamma,\alpha,\beta$: $\Gamma,\alpha\vd\beta$ выполнено тогда и только тогда, когда $\Gamma\vd\alpha\to\beta$.
\end{thm}
\begin{thm}[о полноте исчисления высказываний]
Если $\models\alpha$, то $\vd\alpha$.
\end{thm}

\textbf{Идея доказательства теоремы о дедукции.}
\begin{itemize}
\item ($\Leftarrow$) Если $\delta_1,\dots,\delta_{n-1},\alpha\to\beta$ --- вывод из $\Gamma$, то
  $\delta_1,\dots,\alpha\to\beta,\ \alpha,\ \beta$ --- вывод из $\Gamma,\alpha$ (гипотеза $\alpha$ и M.P.).
\item ($\Rightarrow$) По выводу $\delta_1,\dots,\delta_n$ из $\Gamma,\alpha$ строим вывод $\Gamma\vd\alpha\to\delta_i$ индукцией; новые формулы в «черновик» $\alpha\to\delta_1,\dots,\alpha\to\delta_n$ добавляем по случаям:
  \begin{enumerate}
  \item $\delta_i$ --- аксиома или из $\Gamma$: $\delta_i$, $\ \delta_i\to\alpha\to\delta_i$ (сх.\ 1), $\ \alpha\to\delta_i$ (M.P.).
  \item $\delta_i\equiv\alpha$: вывод $\vd\alpha\to\alpha$ из пяти формул (билет~\hyperref[T:6]{6}).
  \item $\delta_i$ по M.P. из $\delta_j$ и $\delta_k\equiv\delta_j\to\delta_i$: из $\alpha\to\delta_j$ и $\alpha\to\delta_j\to\delta_i$ с помощью сх.\ 2 $(\alpha\to\delta_j)\to(\alpha\to\delta_j\to\delta_i)\to(\alpha\to\delta_i)$ и двух M.P.
  \end{enumerate}
\end{itemize}
Доказательство \emph{конструктивно}: даёт метод перестроения вывода.

\textbf{Полезные леммы:} (а) контрапозиция $\vd(\alpha\to\beta)\to(\neg\beta\to\neg\alpha)$ (из сх.\ 9, 1 и дедукции); (б) закон
исключённого третьего $\vd\alpha\vee\neg\alpha$ (через сх.\ 6, 7, 9, 10 и контрапозицию); (в) об исключении допущения: если $\Gamma,\rho\vd\alpha$ и $\Gamma,\neg\rho\vd\alpha$, то $\Gamma\vd\alpha$ (дедукция, сх.\ 8 и (б)).

\textbf{Идея доказательства полноты (три шага).}
Условное отрицание $\cn{\alpha}=\alpha$, если $\sem{\alpha}=\True$, и $\neg\alpha$, если $\sem{\alpha}=\False$.
\begin{enumerate}
\item \emph{Лемма о связках}: для каждой связки $\star$ доказывается $\cn{\varphi},\cn{\psi}\vd\cn{\varphi\star\psi}$ --- это 14 утверждений-строк таблиц (напр. $\neg\varphi,\neg\psi\vd\neg(\varphi\amp\psi)$, $\varphi,\neg\psi\vd\neg(\varphi\to\psi)$).
\item Пусть $\Xi=\{X_1,\dots,X_n\}$ --- переменные $\alpha$. Тогда $\cn{\Xi}\vd\cn{\alpha}$ при любой оценке (индукция по длине $\alpha$: база --- переменная, переход --- шаг 1).
\item Если $\models\alpha$, то при всех оценках $\cn{\Xi}\vd\alpha$. Индукцией по $n$ убираем гипотезы: из $\cn{X_1..X_n},X_{n+1}\vd\alpha$ и $\cn{X_1..X_n},\neg X_{n+1}\vd\alpha$ по лемме об исключении допущения получаем $\cn{X_1..X_n}\vd\alpha$. В итоге $\vd\alpha$.
\end{enumerate}
Замечание: доказательство конструктивно; вывод экспоненциален по длине (для 3 переменных порядка трёх тысяч строк).

\begin{warn}
\begin{itemize}
\item В теореме о дедукции для \emph{исчисления предикатов} нужны ограничения на правила для кванторов (билеты по лекции 5, вне этого коллоквиума).
\item Вторую часть теоремы о дедукции (о $\Rightarrow$) нельзя получить наивной припиской $\alpha\to$ ко всем формулам: черновик $\alpha\to\delta_i$ не является выводом (пример $\alpha:=A$, $\delta_1:=A\to B\to A$).
\end{itemize}
\end{warn}

% ======================================================================
\group{2. Топологическое пространство}{Билеты 9--13: определение, примеры, открытые и замкнутые множества, связность, непрерывность.\par
\hyperref[T:9]{9}\ \hyperref[T:10]{10}\ \hyperref[T:11]{11}\ \hyperref[T:12]{12}\ \hyperref[T:13]{13}}

\ticket{2(а)}{Топологическое пространство: определение}
\begin{short}
Пара $\langle X,\Omega\rangle$, $\Omega\subseteq\mathcal{P}(X)$: $\varnothing,X\in\Omega$; $\Omega$ замкнуто относительно конечных пересечений и произвольных объединений.
Элементы $\Omega$ --- открытые множества.
\end{short}

\begin{dfn}
\emph{Топологическое пространство} --- упорядоченная пара $\langle X,\Omega\rangle$, где $X$ --- множество (носитель), $\Omega\subseteq\mathcal{P}(X)$ (топология), причём:
\begin{enumerate}
\item $\varnothing\in\Omega$ и $X\in\Omega$;
\item если $A_1,\dots,A_n\in\Omega$, то $A_1\cap\dots\cap A_n\in\Omega$ (конечные пересечения);
\item если $\{A_\alpha\}$ --- произвольное семейство множеств из $\Omega$, то $\bigcup_\alpha A_\alpha\in\Omega$ (произвольные объединения).
\end{enumerate}
Элементы $\Omega$ называют \emph{открытыми} множествами.
\end{dfn}

\textbf{Комментарии.}
\begin{itemize}
\item Пустое объединение даёт $\varnothing$, пустое пересечение --- $X$, так что пункт 1 следовало бы выводить; его включают для наглядности.
\item Условия независимы: можно привести пример $\Omega$, нарушающий ровно одно условие (напр., $\Omega$ без $X$; $\Omega=\{\varnothing,\{1\},\{2\},X\}$ с $X=\{1,2,3\}$ нарушает замкнутость относительно объединения: $\{1\}\cup\{2\}$ отсутствует; $\Omega=\{\varnothing,\{1,2\},\{2,3\},X\}$ нарушает пересечение: $\{2\}\notin\Omega$ ).
\item Для конечного пересечения достаточно парных пересечений (индукция).
\item Топология --- обобщение понятия «близости» без метрики; для логики важно, что открытые множества образуют решётку (билет~\hyperref[T:18]{18}) и интерпретируют интуиционистские связки (билет~\hyperref[T:17]{17}, \hyperref[T:22]{22}).
\end{itemize}

\textbf{Модель интуиционистской логики.} Для пространства $\langle X,\Omega\rangle$ положим $V=\Omega$, «истина» $=X$ и
\[
\sem{\alpha\amp\beta}=\sem{\alpha}\cap\sem{\beta},\ \ \sem{\alpha\vee\beta}=\sem{\alpha}\cup\sem{\beta},\ \
\sem{\neg\alpha}=(\cmpl\sem{\alpha})^\circ,\ \ \sem{\alpha\to\beta}=(\cmpl\sem{\alpha}\cup\sem{\beta})^\circ,
\]
где $\cmpl$ --- дополнение, $^\circ$ --- внутренность. Формула истинна, если её значение равно $X$.

\begin{warn}
\begin{itemize}
\item Объединение \emph{бесконечного} семейства открытых открыто, а пересечение бесконечного --- нет, в общем случае ($\bigcap_n(-\tfrac1n,\tfrac1n)=\{0\}$ в $\bbR$).
\item «Открытое» --- не антоним «замкнутого»: бывают множества и открытые, и замкнутые, и ни те ни другие (билет~\hyperref[T:11]{11}).
\item Бывают пространства, где произвольные пересечения открытых открыты (дискретная, антидискретная, на конечном множестве, Александрова).
\end{itemize}
\end{warn}

% ----------------------------------------------------------------------
\ticket{2(б)}{Примеры топологий (стрелка, Зарисский, деревья и графы, дискретная и антидискретная, евклидова)}
\begin{short}
Дискретная: $\mathcal{P}(X)$. Антидискретная: $\{\varnothing,X\}$. Стрелка: $\{(x,+\infty)\}$. Зарисский: пустое и множества с конечным дополнением. Евклидова: база из интервалов.
Дерево: открыты множества, замкнутые вверх по «потомку».
\end{short}

\begin{enumerate}
\item \textbf{Дискретная} топология на $X$: $\Omega=\mathcal{P}(X)$. Все аксиомы тривиальны; все подмножества открыты (и замкнуты). Для $|X|>1$ несвязна.
\item \textbf{Антидискретная} (тривиальная): $\Omega=\{\varnothing,X\}$. Связна.
\item \textbf{Топология стрелки} на $\bbR$: $\Omega=\{(x,+\infty)\mid x\in\{-\infty\}\cup\bbR\cup\{+\infty\}\}$ (при $x=-\infty$ получаем $\bbR$, при $+\infty$ --- $\varnothing$). Объединение лучей --- луч (с $\sup$ левых концов), пересечение --- луч с большим левым концом.
\item \textbf{Топология Зарисского} на $X$: $\Omega=\{A\subseteq X\mid X\setminus A\text{ конечно}\}\cup\{\varnothing\}$. Пересечение двух множеств с конечными дополнениями имеет дополнение $=$ объединение конечных = конечно; объединение --- дополнение $\subseteq$ конечного. Обратите внимание: на бесконечном $X$ пересечение бесконечного семейства открытых \emph{не} обязано быть открытым: в $\mathbb{N}$ множества $U_n=\mathbb{N}\setminus\{1,\dots,n\}$ открыты, а $\bigcap U_n=\varnothing$ открыто --- но $\bigcap_{a\ \text{чётное}}(\mathbb{N}\setminus\{a\})=\{\text{нечётные}\}$ имеет бесконечное дополнение и не открыто.
\item \textbf{Евклидова} топология на $\bbR$: база $\mathcal{B}=\{(a,b)\mid a,b\in\bbR\}$, т.\,е. $\Omega=\{\bigcup S\mid S\subseteq\mathcal{B}\}$; множество $A$ открыто тогда и только тогда, когда вместе с каждой точкой содержит интервал вокруг неё. (Аналогично на $\bbR^n$ шарами.)
\item \textbf{Топология на дереве (лесе)}. Пусть $V$ --- вершины, $a\preceq b$, если $b$ --- потомок $a$ (или $a=b$). Множество $A\subseteq V$ \emph{открыто}, если из $a\in A$ и $a\preceq b$ следует $b\in A$ (вместе с вершиной --- все её потомки). Проверка: пустое и $V$ открыты; пересечение и объединение множеств, замкнутых вверх по $\preceq$, замкнуты вверх. Для произвольного ориентированного графа ту же конструкцию (открыто --- содержит вместе с точкой всех достижимых) определяют аналогично: получаем топологическое пространство (отношение достижимости --- предпорядок). Для частичного порядка вообще: $\Omega_\preceq=$ все \emph{верхние множества} (топология Александрова); её базу образуют $\{t\mid x\preceq t\}$.
\item \textbf{Топология делимости} на $\mathbb{N}$: открыто множество, содержащее число вместе со всеми кратными (это пример предпорядка $a\preceq b\iff a\mid b$).
\item \textbf{Топология-подпространство}: для $A\subseteq X$ положим $\Omega_A=\{S\cap A\mid S\in\Omega\}$ (билет~\hyperref[T:12]{12}).
\end{enumerate}

\textbf{Рисунок: дерево $W_1\to W_2,W_3,W_4$; $W_2\to W_5$; $W_4\to W_6$.} Открытые: $\{W_2,W_5,W_3\}$, $\{W_4,W_6\}$ и всё, что замкнуто вверх. Не открыты: $\{W_1,W_3\}$ (нет потомков $W_2,W_4$), $\{W_2\}$ (нет потомка $W_5$).

\textbf{Примеры, нарушающие аксиомы.} (на $X=\{1,2,3\}$)
\begin{itemize}
\item $\Omega=\{\varnothing,\{1\},\{2\},X\}$ --- нет $\{1,2\}=\{1\}\cup\{2\}$;
\item $\Omega=\{\varnothing,\{1,2\},\{2,3\},X\}$ --- нет $\{2\}=\{1,2\}\cap\{2,3\}$;
\item $\Omega=\{\varnothing,\{1\}\}$ --- нет $X$.
\end{itemize}

\begin{warn}
\begin{itemize}
\item Для стрелки: открытые множества --- именно лучи $(x,+\infty)$ (не $[x,+\infty)$ и не интервалы). Иногда (в источниках) встречается запись $[a,+\infty)$ --- следите за определением, принятым на лекции.
\item В топологии Зарисского на $\bbR$ любые два непустых открытых множества пересекаются, поэтому пространство связно и «не хаусдорфово».
\item Евклидова топология на $\bbR$ \emph{не} дискретна: $\{0\}$ не открыто.
\end{itemize}
\end{warn}

% ----------------------------------------------------------------------
\ticket{2(в)}{Открытые и замкнутые множества. Внутренность, замыкание}
\begin{short}
Открытое --- элемент $\Omega$; замкнутое --- дополнение открытого. Внутренность $A^\circ$ --- наибольшее открытое $\subseteq A$ (объединение всех открытых подмножеств); замыкание $\overline{A}$ --- наименьшее замкнутое $\supseteq A$; $\overline{A}=\cmpl((\cmpl A)^\circ)$.
\end{short}

\begin{dfn}
$A\subseteq X$ \emph{открыто}, если $A\in\Omega$; \emph{замкнуто}, если $X\setminus A\in\Omega$. Множество может быть открытым и замкнутым одновременно (\emph{открыто-замкнутое}), либо ни тем, ни другим.
\end{dfn}
\begin{dfn}
\emph{Внутренность} $A^\circ$ --- наибольшее по включению открытое множество, содержащееся в $A$. \emph{Замыкание} $\overline{A}$ --- наименьшее по включению замкнутое множество, содержащее $A$.
\end{dfn}
\begin{thm}
$A^\circ$ определена для любого $A$: $A^\circ=\bigcup\{Q\in\Omega\mid Q\subseteq A\}$.
\end{thm}
\begin{proof}
Пусть $V=\{Q\in\Omega\mid Q\subseteq A\}$. $\bigcup V\in\Omega$ как объединение открытых и $\bigcup V\subseteq A$; любое $Q\in V$ содержится в $\bigcup V$. Значит $\bigcup V$ --- наибольший элемент $V$.
\end{proof}
Двойственно, $\overline{A}=\bigcap\{F\mid F\text{ замкнуто},A\subseteq F\}$, и $\overline{A}=X\setminus(X\setminus A)^\circ$.

\textbf{Свойства.} $A^\circ\subseteq A\subseteq\overline{A}$; $A$ открыто $\iff A=A^\circ$; $A$ замкнуто $\iff A=\overline{A}$; $(A\cap B)^\circ=A^\circ\cap B^\circ$; $(A^\circ)^\circ=A^\circ$; $\overline{A\cup B}=\overline{A}\cup\overline{B}$.

\textbf{Примеры.} В $\bbR$: $A=(0,1]$: $A^\circ=(0,1)$, $\overline{A}=[0,1]$. $\mathbb{Q}^\circ=\varnothing$, $\overline{\mathbb{Q}}=\bbR$. В стрелке: $[0,1]^\circ=\varnothing$.
Открыто-замкнутые в $\bbR$ с евклидовой топологией: только $\varnothing$ и $\bbR$ (связность, билет~\hyperref[T:12]{12}). В дискретной топологии открыто-замкнуто всё. Можно построить топологию на $\bbR$ ровно с четырьмя открыто-замкнутыми множествами (например, $\bbR=(-\infty,0)\sqcup[0,\infty)$ с евклидовой на каждой части) и вообще с $2^k$; ровно \emph{три} невозможно: открыто-замкнутые множества образуют булеву алгебру, а в ней $2^k$ элементов или бесконечно.

\textbf{В логике.} Внутренность определяет интуиционистскую импликацию и отрицание: $\sem{\alpha\to\beta}=(\cmpl\sem{\alpha}\cup\sem{\beta})^\circ$, $\sem{\neg\alpha}=(\cmpl\sem{\alpha})^\circ$.
Пример: $A=(0,1)\subset\bbR$: $\sem{A\vee\neg A}=(0,1)\cup(-\infty,0)\cup(1,\infty)=\bbR\setminus\{0,1\}\neq\bbR$, поэтому $A\vee\neg A$ не общезначима.

\begin{warn}
\begin{itemize}
\item $(A\cup B)^\circ\ne A^\circ\cup B^\circ$ в общем случае: $A=[0,1],B=[1,2]$: $(A\cup B)^\circ=(0,2)$, а $A^\circ\cup B^\circ=(0,1)\cup(1,2)$.
\item Внутренность через решёточный язык: $A^\circ=\inf_\Omega\{A\}$ (наибольшая нижняя грань в решётке $\Omega$).
\end{itemize}
\end{warn}

% ----------------------------------------------------------------------
\ticket{2(г)}{Связность, линейная связность}
\begin{short}
Пространство связно, если его нельзя разбить на два непустых непересекающихся открытых множества (эквивалентно: открыто-замкнуты только $\varnothing$ и $X$).
Линейно связно: любые две точки соединяются путём. Линейно связное $\Rightarrow$ связное; обратное неверно.
\end{short}

\begin{dfn}
$\langle X,\Omega\rangle$ \emph{связно}, если не существует $A,B\in\Omega$ с $A\cup B=X$, $A\cap B=\varnothing$, $A\ne\varnothing$, $B\ne\varnothing$.
Эквивалентно: единственные открыто-замкнутые множества --- $\varnothing$ и $X$ (если $A$ открыто-замкнуто, то $B=X\setminus A$ тоже открыто).
\end{dfn}
\begin{dfn}[подпространство]
$\langle X_1,\Omega_1\rangle$ --- подпространство $\langle X,\Omega\rangle$, если $X_1\subseteq X$ и $\Omega_1=\{S\cap X_1\mid S\in\Omega\}$. Пример: $[0,1)$ открыто в $[0,1]$, так как $[0,1)=(-1,1)\cap[0,1]$. Множество $A\subseteq X$ \emph{связно}, если связно как подпространство.
\end{dfn}
\begin{dfn}[путь, линейная связность]
\emph{Путь} между $a$ и $b$ --- непрерывная $f:[0,1]\to X$, $f(0)=a$, $f(1)=b$. Множество \emph{линейно связно}, если любые две его точки соединяются путём, лежащим в множестве.
\end{dfn}

\textbf{Примеры.}
\begin{itemize}
\item $(0,1]\cup[2,3)\subset\bbR$ несвязно: $A=(0,1]$, $B=[2,3)$ открыты в подпространстве.
\item Дискретное пространство с $|X|>1$ несвязно: $A=\{a\}$, $B=X\setminus A$.
\item $\mathbb{Q}\subset\bbR$ несвязно (разбиваем иррациональной точкой), но топология на $\mathbb{Q}$ \emph{не} дискретна (нет изолированных точек).
\item Антидискретное пространство и пространство с топологией Зарисского на бесконечном множестве связны.
\end{itemize}

\begin{lmm}[деревья]
Лес связен (является одним деревом) $\iff$ соответствующее топологическое пространство связно.
\end{lmm}
\begin{proof}
($\Rightarrow$) Пусть дерево с корнем $v$ и $A,B$ открыты, непусты, дизъюнктны, $A\cup B=V$; пусть $v\in A$. Тогда $A\supseteq\{x\mid v\preceq x\}=V$, $B=\varnothing$ --- противоречие.
($\Leftarrow$) Если корней $v_1,\dots,v_k$ несколько ($k\ge2$), множества $A_i=\{x\mid v_i\preceq x\}$ открыты, непусты, дизъюнктны и $\bigcup A_i=V$ --- несвязно.
\end{proof}
Для частичного порядка из конечного ациклического графа: слабая связность графа $\iff$ связность топологического пространства порядка.

\begin{thm}
Линейно связное множество связно.
\end{thm}
\begin{proof}
Пусть $X=A\sqcup B$ (открытые, непустые); возьмём $a\in A$, $b\in B$ и путь $f:[0,1]\to X$ от $a$ к $b$. Тогда $f^{-1}(A)$, $f^{-1}(B)$ --- открытые непустые непересекающиеся, покрывающие $[0,1]$; но $[0,1]$ связен (контрадикция). (Связность $[0,1]$: если $[0,1]=U\sqcup W$, $0\in U$, положим $s=\sup\{t\mid[0,t]\subseteq U\}$ и разберём $s\in U$ или $s\in W$.)
\end{proof}
Обратное неверно: «топологическая синусоида» $\{(x,\sin\frac1x)\mid x>0\}\cup\{(0,y)\mid|y|\le1\}$ связна, но не линейно связна.

\textbf{Примеры путей.} В дискретном $\mathbb{N}$ путь постоянен (иначе $[0,1]$ распалось бы на прообразы точек). В топологии стрелки путь --- функция $f:[0,1]\to\bbR$, для которой $f^{-1}((x,+\infty))$ открыт для всех $x$ (полунепрерывная снизу).

\begin{warn}
\begin{itemize}
\item Связность множества --- свойство его \emph{подпространства}: разбивать надо множествами $U\cap A$, $V\cap A$ с открытыми $U,V\subseteq X$, а не открытыми подмножествами самого $A$ в объемлющей топологии.
\item Не путайте связность в топологии и связность графа: для конечного графа слабая связность эквивалентна топологической (задача ДЗ); для бесконечных нужно разбирать.
\end{itemize}
\end{warn}

% ----------------------------------------------------------------------
\ticket{2(д)}{Непрерывные функции}
\begin{short}
$f:X\to Y$ непрерывна, если прообраз любого открытого множества открыт. Композиция непрерывных непрерывна. Путь --- непрерывная $[0,1]\to X$.
\end{short}

\begin{dfn}
Функция $f:\langle X,\Omega_X\rangle\to\langle Y,\Omega_Y\rangle$ \emph{непрерывна}, если для любого $U\in\Omega_Y$ множество $f^{-1}(U)$ лежит в $\Omega_X$.
\end{dfn}

\textbf{Примеры.}
\begin{itemize}
\item Любая $f:\mathbb{N}\to\bbR$ при дискретной топологии на $\mathbb{N}$ непрерывна (любое множество в $\mathbb{N}$ открыто).
\item Постоянная функция всегда непрерывна; тождественная $\mathrm{id}:\langle X,\Omega\rangle\to\langle X,\Omega'\rangle$ непрерывна тогда и только тогда, когда $\Omega'\subseteq\Omega$.
\item В евклидовой топологии на $\bbR$ определение совпадает с $\varepsilon$--$\delta$.
\item Для конечных деревьев/порядков (топология Александрова) непрерывные функции --- монотонные (сохраняющие $\preceq$).
\end{itemize}

\textbf{Свойства.} Композиция непрерывных непрерывна. Образ связного пространства при непрерывной функции связен. Образ компактного компактен.

\begin{dfn}[ещё два понятия топологии]
\emph{Компактное} множество --- из любого его открытого покрытия можно выбрать конечное подпокрытие. \emph{Путь} --- непрерывное отображение отрезка $[0,1]$ в пространство (билет~\hyperref[T:12]{12}).
\end{dfn}

\textbf{Двойственное определение.} «Новая непрерывность» --- образ открытого открыт (\emph{открытое отображение}) --- от непрерывности не следует и её не влечёт: $x\mapsto x^2$ непрерывна, но $(-1,1)\mapsto[0,1)$ не открыто; проекция $\bbR^2\to\bbR$ открыта; постоянная непрерывна и не открыта (если $|Y|>1$).

\begin{warn}
\begin{itemize}
\item Условие «прообраз открытого открыт», а не «образ»: частая ошибка.
\item Непрерывность зависит от выбора топологий на обоих пространствах.
\end{itemize}
\end{warn}

% ======================================================================
\group{3. Гильбертов вывод и натуральный вывод}{Билет 14.\par\hyperref[T:14]{14}}

\ticket{3}{Гильбертов вывод и натуральный вывод}
\begin{short}
Гильбертов вывод: много аксиом, одно правило (M.P.), контекст --- метаязык. Натуральный вывод: мало аксиом, много правил введения и удаления, секвенция $\Gamma\vd\alpha$ --- часть языка, вывод --- дерево. Системы эквивалентны.
\end{short}

\textbf{Гильбертов стиль.} Язык: $\neg,\amp,\vee,\to$. Аксиомы --- 10 схем (билет~\hyperref[T:6]{6}); правило вывода --- Modus Ponens. Вывод --- линейная последовательность формул; гипотезы $\Gamma$ --- внешние по отношению к формулам.
\emph{Достоинство}: просто определить; \emph{недостаток}: выводы длинные и нечитаемые (для $A\to A$ уже 5 строк), без теоремы о дедукции работать неудобно.

\textbf{Натуральный (естественный) вывод для ИИВ.}
Язык: $\amp,\vee,\to,\bot$ ($\neg\alpha:=\alpha\to\bot$). Формулы вывода --- \emph{секвенции} $\Gamma\vd\alpha$ (контекст теперь \emph{часть предметного языка}; поэтому $\vd$ внутри и снаружи --- разные символы, различают $|\alpha|_{\text{г}}$ и $|\alpha|_{\text{н}}$).
Вывод записывается деревом: под чертой заключение, над чертой посылки.
\begin{align*}
&\text{Аксиома:}&&\infer[\text{(акс.)}]{\Gamma,\alpha\vd\alpha}{}\\
&\text{Введение связок:}&&\infer{\Gamma\vd\alpha\to\beta}{\Gamma,\alpha\vd\beta}\quad
\infer{\Gamma\vd\alpha\vee\beta}{\Gamma\vd\alpha}\quad\infer{\Gamma\vd\alpha\vee\beta}{\Gamma\vd\beta}\quad
\infer{\Gamma\vd\alpha\amp\beta}{\Gamma\vd\alpha&\Gamma\vd\beta}\\
&\text{Удаление связок:}&&\infer{\Gamma\vd\beta}{\Gamma\vd\alpha&\Gamma\vd\alpha\to\beta}\quad
\infer{\Gamma\vd\gamma}{\Gamma\vd\alpha\to\gamma&\Gamma\vd\beta\to\gamma&\Gamma\vd\alpha\vee\beta}\\
&&&\infer{\Gamma\vd\alpha}{\Gamma\vd\alpha\amp\beta}\quad\infer{\Gamma\vd\beta}{\Gamma\vd\alpha\amp\beta}\quad
\infer{\Gamma\vd\alpha}{\Gamma\vd\bot}\ \text{(принцип взрыва)}
\end{align*}
Пример: $A\amp B\vd B\amp A$:
\[
\infer[(\text{введ.}\amp)]{A\amp B\vd B\amp A}{
 \infer[(\text{удал.}\amp)]{A\amp B\vd B}{\infer[(\text{акс.})]{A\amp B\vd A\amp B}{}}
 &
 \infer[(\text{удал.}\amp)]{A\amp B\vd A}{\infer[(\text{акс.})]{A\amp B\vd A\amp B}{}}}
\]

\textbf{Классическая логика.} Принцип взрыва $\bot\vd\alpha$ заменяют \emph{снятием двойного отрицания}:
\[
\infer{\Gamma\vd\alpha}{\Gamma\vd(\alpha\to\bot)\to\bot}
\]
Таким образом классический и интуиционистский варианты различаются одним правилом.

\textbf{Эквивалентность стилей.} По индукции можно показать: если $\Gamma\vd_{\text{ги}}\alpha$, то $|\Gamma|_{\text{н}}\vd_{\text{ни}}|\alpha|_{\text{н}}$ (гильбертовскую аксиоматику выводим в натуральной), и наоборот $\Gamma\vd_{\text{ни}}\alpha$ даёт $|\Gamma|_{\text{ги}}\vd_{\text{г}}|\alpha|_{\text{ги}}$ (по теореме о дедукции правило введения $\to$ --- это сама теорема о дедукции). Индексы снизу: $\vd_{\text{и}}$, $\vd_{\text{к}}$ --- метаязык.

\textbf{Различия.}
\begin{center}
\begin{tabular}{p{4cm}|p{5.2cm}|p{5.2cm}}
&Гильбертов&Натуральный\\\hline
Аксиомы/правила&10 схем, M.P.&1 аксиома, правила введения/удаления\\
Контекст&метаязык&в секвенции\\
Вывод&последовательность&дерево\\
Ложь&через $\neg$ (сх. 9, 10/10и)&константа $\bot$\\
Удобство&простое определение&естественные рассуждения, соответствие Карри--Ховарду
\end{tabular}
\end{center}

\begin{warn}
\begin{itemize}
\item В натуральном выводе есть \emph{ровно одно} правило на каждое введение и удаление; правило удаления дизъюнкции --- аналог сх. 8.
\item Дерево вывода соответствует программе (билет~\hyperref[T:28]{28}): введение $\to$ --- $\lambda$-абстракция, удаление --- применение.
\end{itemize}
\end{warn}

% ======================================================================
\group{4. Интуиционистское исчисление высказываний}{Билеты 15--27: от философии конструктивизма до разрешимости.\par
\hyperref[T:15]{15}\ \hyperref[T:16]{16}\ \hyperref[T:17]{17}\ \hyperref[T:18]{18}\ \hyperref[T:19]{19}\ \hyperref[T:20]{20}\ \hyperref[T:21]{21}\ \hyperref[T:22]{22}\ \hyperref[T:23]{23}\ \hyperref[T:24]{24}\ \hyperref[T:25]{25}\ \hyperref[T:26]{26}\ \hyperref[T:27]{27}}

\ticket{4(а)}{Доказательства чистого существования}
\begin{short}
Доказательство существования объекта без его предъявления (обычно от противного или через закон исключённого третьего). Примеры: теорема Брауэра о неподвижной точке; <<иррациональное в иррациональной степени>>.
\end{short}

\textbf{Идея.} Доказательство \emph{чистого существования} утверждает $\exists x\,P(x)$, не указывая никакого $x$ (и никакого способа его найти). В конструктивной математике такое доказательство не принимается --- объект должен быть построен.

\textbf{Пример 1 (Брауэр).} Любое непрерывное отображение $f$ шара в $\bbR^n$ в себя имеет неподвижную точку.
Доказательство: пусть нет неподвижных точек; тогда для каждой $X$ луч от $f(X)$ через $X$ выходит на границу шара --- получаем непрерывное отображение шара на границу, оставляющее границу на месте (ретракция), а такого не существует (теорема без доказательства). Противоречие. Метод построения неподвижной точки не предложен --- это типичное доказательство чистого существования (от противного).

\textbf{Пример 2.} \emph{Существуют иррациональные $a,b$, что $a^b\in\mathbb{Q}$.}
Рассмотрим $\sqrt2^{\sqrt2}$.
\begin{enumerate}
\item Если оно рационально, то $a=b=\sqrt2$.
\item Если иррационально, то $a=\sqrt2^{\sqrt2}$, $b=\sqrt2$: $a^b=\left(\sqrt2^{\sqrt2}\right)^{\sqrt2}=\sqrt2^{\,2}=2\in\mathbb{Q}$.
\end{enumerate}
Доказательство \emph{неконструктивно}: оно не говорит, какой из двух случаев имеет место. Причина --- использован закон исключённого третьего $P\vee\neg P$ для $P=$ «$\sqrt2^{\sqrt2}\in\mathbb{Q}$». Формально: доказаны $P\to S$ и $\neg P\to S$; чтобы получить $S$, нужно $P\vee\neg P$; в BHK-интерпретации дизъюнкция требует знать, какая из альтернатив верна, поэтому рассуждение не является доказательством в этом смысле. (На самом деле $\sqrt2^{\sqrt2}$ иррационально --- теорема Гельфонда--Шнайдера; конструктивный пример: $a=\sqrt2$, $b=2\log_2 3$, $a^b=3$.)

\textbf{Программистский смысл.} Такая теорема не позволяет написать программу, ищущую пару чисел. Конструктивное доказательство, напротив, \emph{является} программой (билет~\hyperref[T:28]{28}).

\begin{warn}
\begin{itemize}
\item Доказательство от противного для утверждений вида $\neg\alpha$ конструктивно приемлемо ($\neg\alpha=\alpha\to\bot$); неприемлемо снятие двойного отрицания $\neg\neg\alpha\to\alpha$.
\item Чистое существование не «неверно» --- оно просто слабее по информации и не принимается интуиционизмом.
\end{itemize}
\end{warn}

% ----------------------------------------------------------------------
\ticket{4(б)}{Интуиционизм (программа Брауэра). BHK-интерпретация (в том числе расшифровка названия)}
\begin{short}
Брауэр: математика --- свободная деятельность ума, логика вторична. BHK = Брауэр, Гейтинг, Колмогоров. Истина = наличие построения: конъюнкция --- пара, дизъюнкция --- выбор с меткой, импликация --- функция, $\bot$ не имеет построения, $\neg\alpha=\alpha\to\bot$.
\end{short}

\textbf{Интуиционизм} (Брауэр, <<Over de Grondslagen der Wiskunde>>, 1907). Основные положения:
\begin{enumerate}
\item Математика не формальна.
\item Математика независима от окружающего мира.
\item Математика не зависит от логики --- это логика зависит от математики.
\end{enumerate}
Математик не рассуждает в стиле классической логики; логика должна описывать «интуитивный» стиль. Побочно: исключаются «неконструктивные» объекты (в частности, множество $\{x\mid x\notin x\}$ предъявить нельзя).

\textbf{BHK-интерпретация} (по фамилиям: \textbf{Б}рауэр --- \textbf{Г}ейтинг (Heyting) --- \textbf{К}олмогоров). Индуктивно определяется, что значит «$\alpha$ построено/доказано»:
\begin{itemize}
\item $\alpha\amp\beta$ построено, если построены $\alpha$ и $\beta$;
\item $\alpha\vee\beta$ построено, если построено $\alpha$ или построено $\beta$, и мы знаем, что именно;
\item $\alpha\to\beta$ построено, если есть способ преобразовать любое построение $\alpha$ в построение $\beta$;
\item $\bot$ --- конструкция, не имеющая построения;
\item $\neg\alpha$ построено, если построено $\alpha\to\bot$ (из построения $\alpha$ получаем противоречие).
\end{itemize}

\textbf{Следствия.}
\begin{itemize}
\item $\alpha\vee\neg\alpha$ в общем случае не имеет построения: для нерешённой проблемы ($P=NP$?) нет ни построения $\alpha$, ни $\neg\alpha$.
\item Формула $(A\to B)\vee(B\to C)\vee(C\to A)$ классически общезначима, но в BHK неприемлема: «материальная» импликация (надо посмотреть в окно) против «формальной» (причинно-следственной связи нет).
\end{itemize}

\textbf{Формализация} --- не первична: «никакая формализация не является первичной; мы выбираем те или иные правила, потому что они соответствуют целям, и вольны их менять». Наиболее распространённая: гильбертовская аксиоматика КИВ с заменой аксиомы (10) $\neg\neg\alpha\to\alpha$ на (10и) $\alpha\to\neg\alpha\to\beta$.

\textbf{Проверка:} в интуиционистском исчислении доказуемо $\alpha\to\neg\neg\alpha$ и $\neg\neg\neg\alpha\to\neg\alpha$, но недоказуемо $\neg\neg\alpha\to\alpha$ и $\alpha\vee\neg\alpha$ (билеты~\hyperref[T:17]{17}, \hyperref[T:24]{24}).

\begin{warn}
\begin{itemize}
\item В BHK-интерпретации нет формальной модели --- модели появляются позже (топологическая, Крипке, решётки).
\item Название: Hey\textbf{t}ing = Гейтинг, поэтому «алгебра Гейтинга» и «Heyting algebra» --- один и тот же человек.
\end{itemize}
\end{warn}

% ----------------------------------------------------------------------
\ticket{4(в)}{Связь с формальными логиками. Закон исключённого третьего, принцип взрыва; классическая, интуиционистская, позитивная логики}
\begin{short}
Позитивная $\subset$ минимальная $\subset$ интуиционистская $\subset$ классическая. Различия: аксиомы 9 (отрицание), «взрыв» $\alpha\to\neg\alpha\to\beta$ и снятие двойного отрицания $\neg\neg\alpha\to\alpha$. В ИИВ $\alpha\vee\neg\alpha$ недоказуем.
\end{short}

\textbf{Аксиоматики.}
\begin{itemize}
\item \emph{Позитивная логика}: аксиомы 1--8 (без отрицания) --- конъюнкция, дизъюнкция, импликация.
\item \emph{Минимальная}: добавлена 9 $(\alpha\to\beta)\to(\alpha\to\neg\beta)\to\neg\alpha$ (опровержение допущения, приводящего к противоречию).
\item \emph{Интуиционистская (ИИВ)}: минимальная + (10и) $\alpha\to\neg\alpha\to\beta$ (принцип взрыва: из противоречия следует всё).
\item \emph{Классическая (КИВ)}: минимальная + (10) $\neg\neg\alpha\to\alpha$ (снятие двойного отрицания); из (10) принцип взрыва выводим, поэтому КИВ сильнее ИИВ.
\end{itemize}
Каждая следующая логика доказывает больше. В КИВ доказуемы тавтологии, в ИИВ --- меньше.

\textbf{Закон исключённого третьего (ЗИТ)} $\alpha\vee\neg\alpha$: доказуем в КИВ, не доказуем в ИИВ (контрмодель --- билеты~\hyperref[T:22]{22}, \hyperref[T:24]{24}: топология $\bbR$, $A=(0,1)$ или крипкевская модель из трёх миров).

\textbf{Принцип взрыва (ex falso)} $\alpha\to\neg\alpha\to\beta$ (иначе: $\bot\vd\beta$). Принимается в ИИВ и КИВ, не принимается в минимальной логике.

\textbf{Что доказуемо в ИИВ:} $\alpha\to\neg\neg\alpha$, $\neg\neg\neg\alpha\to\neg\alpha$, $(\alpha\to\beta)\to(\neg\beta\to\neg\alpha)$, $\neg(\alpha\amp\neg\alpha)$, $\neg\neg(\alpha\vee\neg\alpha)$ (двойное отрицание ЗИТ!).
\textbf{Недоказуемо в ИИВ:} $\neg\neg\alpha\to\alpha$, $\alpha\vee\neg\alpha$, закон Пирса $((\alpha\to\beta)\to\alpha)\to\alpha$, $(\alpha\to\beta)\to\neg\alpha\vee\beta$, $\neg(\alpha\amp\beta)\to\neg\alpha\vee\neg\beta$.

\textbf{Кольцо эквивалентностей (над минимальной логикой, ДЗ 4).}
\[
\neg\neg\alpha\to\alpha\ \Longrightarrow\ ((\alpha\to\beta)\to\alpha)\to\alpha\ \Longrightarrow\ \alpha\vee\neg\alpha\ \Longrightarrow\ \neg\neg\alpha\to\alpha\quad\text{(последний шаг --- при принципе взрыва).}
\]
Закон Пирса и принцип взрыва независимы; из закона Пирса не следует снятие двойного отрицания (без взрыва).

\textbf{Дизъюнктивное свойство.} В ИИВ $\vd\alpha\vee\beta$ влечёт $\vd\alpha$ или $\vd\beta$ (билет~\hyperref[T:31]{31}) --- согласуется с BHK; в КИВ неверно: $\vd A\vee\neg A$, но ни $A$, ни $\neg A$ не доказуемы.

\textbf{Модели.} КИВ $\leftrightarrow$ булевы алгебры; ИИВ $\leftrightarrow$ алгебры Гейтинга $\leftrightarrow$ топологические пространства $\leftrightarrow$ модели Крипке.

\begin{warn}
\begin{itemize}
\item Не говорите «ИИВ не доказывает $A\vee\neg A$, значит доказывает $\neg(A\vee\neg A)$»: наоборот, $\vd_{\text{и}}\neg\neg(A\vee\neg A)$, т.\,е. $\neg(A\vee\neg A)$ противоречиво --- ЗИТ нельзя опровергнуть.
\item «Позитивная» --- не «положительная»: это фрагмент без отрицания (и без $\bot$).
\end{itemize}
\end{warn}

% ----------------------------------------------------------------------
\ticket{4(г)}{Решётки. Топология как решётка}
\begin{short}
Решётка --- частичный порядок, где у любых двух элементов есть $\sup$ ($+$) и $\inf$ ($\cdot$). Открытые множества с $\subseteq$ образуют решётку ($+=\cup$, $\cdot=\cap$).
\end{short}

\textbf{Предварительные понятия.} Пусть $\langle\mathcal{U},\preceq\rangle$ --- частичный порядок, $X\subseteq\mathcal{U}$.
\begin{itemize}
\item $\mathrm{lwb}_{\mathcal U}X$ --- нижние грани, $\mathrm{upb}_{\mathcal U}X$ --- верхние грани;
\item \emph{минимальный} $m\in X$: нет меньшего ($y\preceq m\Rightarrow y=m$); \emph{максимальный}: нет большего; \emph{наименьший}: $m\preceq y$ при всех $y\in X$; \emph{наибольший};
\item $\inf X=\mathrm{наиб}(\mathrm{lwb}X)$, $\sup X=\mathrm{наим}(\mathrm{upb}X)$.
\end{itemize}
Наименьший/наибольший единственны, минимальных/максимальных может быть много. Пример (делимость на $\mathbb{N}$, $a\preceq b\iff a\mid b$): множество $\{2,3,6\}$ имеет минимальные $2,3$, наименьшего нет, $\inf=1$. Пример: $\{1,1.4,1.41,\dots\}\subset\mathbb{Q}$ не имеет $\sup$ в $\mathbb{Q}$.

\begin{dfn}
\emph{Решётка} --- пара $\langle X,\preceq\rangle$, где $\preceq$ --- частичный порядок и для любых $a,b\in X$ определены $a+b=\sup\{a,b\}$ и $a\cdot b=\inf\{a,b\}$.
\end{dfn}
Примеры: $\langle\mathcal{P}(X),\subseteq\rangle$ --- решётка ($+=\cup$, $\cdot=\cap$); $\langle\mathbb{N},\mid\rangle$ --- решётка ($+=\mathrm{НОК}$, $\cdot=\mathrm{НОД}$); $\langle\mathbb{N}\setminus\{1\},\mid\rangle$ --- \emph{не} решётка (у $\{2,3\}$ нет $\inf$).

\textbf{Топология как решётка.} Для $\langle X,\Omega\rangle$ порядок $\subseteq$ на $\Omega$ делает $\Omega$ решёткой: $A+B=A\cup B$ (объединение открытых открыто и является наименьшей верхней гранью), $A\cdot B=A\cap B$. При этом внутренность определяется решёточно:
\[
A^\circ=\inf{}_\Omega\{A\}=\bigcup\{Q\in\Omega\mid Q\subseteq A\}.
\]
Решётка $\Omega$ имеет $0=\varnothing$ и $1=X$. Это не булева алгебра в общем случае, но (билет~\hyperref[T:21]{21}) псевдобулева: $a\to b=((\cmpl a)\cup b)^\circ$.

\textbf{Решётки и исчисление высказываний.} Если исчисление оценивается значениями из решётки, оценка называется \emph{согласованной}, если $\sem{\alpha\amp\beta}=\sem{\alpha}\cdot\sem{\beta}$, $\sem{\alpha\vee\beta}=\sem{\alpha}+\sem{\beta}$, $\sem{\alpha\to\beta}=\sem{\alpha}\to\sem{\beta}$, $\sem{\neg\alpha}=\ \sim\sem{\alpha}$, $\sem{A\amp\neg A}=0$, $\sem{A\to A}=1$. Любая псевдобулева алгебра, являющаяся согласованной оценкой ИИВ, --- корректная модель: $\vd\alpha\Rightarrow\sem{\alpha}=1$.

\begin{warn}
\begin{itemize}
\item Решётка --- это именно $\sup$ и $\inf$ \emph{любых двух} элементов, а не просто частичный порядок.
\item «Супремум $\ne$ максимум»: максимум лежит в множестве, супремум --- нет, если не принадлежит.
\end{itemize}
\end{warn}

% ----------------------------------------------------------------------
\ticket{4(д)}{Дистрибутивная решётка}
\begin{short}
Дистрибутивная: $a\cdot(b+c)=a\cdot b+a\cdot c$. Любая импликативная решётка дистрибутивна; конечная дистрибутивная --- импликативна. Верно и двойственное $a+b\cdot c=(a+b)(a+c)$.
\end{short}

\begin{dfn}
Решётка \emph{дистрибутивна}, если для любых $a,b,c$: $a\cdot(b+c)=a\cdot b+a\cdot c$.
\end{dfn}
В любой решётке всегда выполнено $a\cdot b+a\cdot c\preceq a\cdot(b+c)$; дистрибутивность --- обратное включение.

\begin{lmm}
В дистрибутивной решётке также $a+b\cdot c=(a+b)\cdot(a+c)$ (двойственный закон).
\end{lmm}
\begin{proof}
$(a+b)(a+c)=(a+b)a+(a+b)c=a+bc+ac$ (закон поглощения $a(a+b)=a$ и дистрибутивность) $=a+bc$, так как $ac\preceq a$.
\end{proof}

\textbf{Основные факты.}
\begin{itemize}
\item \emph{Законы поглощения} верны в любой решётке: $a\cdot(a+b)=a$, $a+a\cdot b=a$.
\item \emph{Монотонность}: $a\preceq b$, $c\preceq d$ $\Rightarrow$ $a+c\preceq b+d$, $a\cdot c\preceq b\cdot d$.
\item \emph{Любая импликативная решётка дистрибутивна.} Доказательство: достаточно $a(b+c)\preceq ab+ac$. Покажем $b\preceq a\to(ab+ac)$ (так как $a\cdot b\preceq ab+ac$) и аналогично $c$; тогда $b+c\preceq a\to(ab+ac)$, откуда $a(b+c)\preceq ab+ac$ по определению $\to$.
\item \emph{Конечная дистрибутивная решётка импликативна}: $a\to b=\sum\{x\mid a\cdot x\preceq b\}$ (сумма конечна; $a\cdot\sum x_i=\sum a x_i\preceq b$).
\item \emph{Критерий} (билет~\hyperref[T:20]{20}): решётка дистрибутивна $\iff$ в ней нет подрешёток-диаманта $M_3$ и пентагона $N_5$.
\end{itemize}

\textbf{Примеры.} $\mathcal{P}(X)$, $\Omega$, цепи (линейно упорядоченные), $\langle\mathbb{N},\mid\rangle$ --- дистрибутивны. Диамант (три атома между $0$ и $1$) и пентагон --- нет.

\textbf{Дистрибутивная, но не импликативная:} $L=\{\text{конечные подмножества }\mathbb{N}\}\cup\{\mathbb{N}\}$ с $\subseteq$ --- подрешётка $\mathcal{P}(\mathbb{N})$, поэтому дистрибутивна, есть $0=\varnothing$ и $1=\mathbb{N}$. Но для $a=\{0\}$, $b=\varnothing$ множество $\{x\in L\mid a\cap x\subseteq b\}$ --- все конечные подмножества $\mathbb{N}\setminus\{0\}$ --- не имеет наибольшего, т.\,е. $a\to b$ не существует.

\begin{warn}
\begin{itemize}
\item Дистрибутивность --- условие на $(+,\cdot)$; псевдодополнение --- дополнительная структура. Из первого не следует второе (в бесконечных решётках).
\item Подрешётка дистрибутивной решётки дистрибутивна (в отличие от импликативности: псевдодополнения могут не сохраниться; билет~\hyperref[T:20]{20}).
\end{itemize}
\end{warn}

% ----------------------------------------------------------------------
\ticket{4(е)}{Подрешётка. Диамант и пентагон; дистрибутивная решётка как решётка без диамантов/пентагонов}
\begin{short}
Подрешётка: подмножество, замкнутое относительно \emph{тех же} $+$ и $\cdot$. Диамант $M_3$ и пентагон $N_5$ недистрибутивны; решётка дистрибутивна $\iff$ не содержит подрешёток, изоморфных $M_3$ и $N_5$.
\end{short}

\begin{dfn}
$\mathcal{L}'=\langle L',\preceq'\rangle$ --- \emph{подрешётка} $\mathcal{L}=\langle L,\preceq\rangle$, если $L'\subseteq L$, $\preceq'$ --- ограничение $\preceq$, и для $a,b\in L'$: $a+_{\mathcal{L}'}b=a+_{\mathcal{L}}b$, $a\cdot_{\mathcal{L}'}b=a\cdot_{\mathcal{L}}b$.
\end{dfn}
Условие совпадения $+$ и $\cdot$ существенно, если ослабить условие на порядок до $(\preceq')\subseteq(\preceq)$: диамант $L'=\{a,b,c,d,e\}$ (с $b+c=e$, $b\cdot c=a$) имеет порядок, содержащийся в порядке решётки $L$, где $c\preceq b$; но $b+_Lc=b\ne e=b+_{L'}c$ и $b\cdot_Lc=c\ne a$, поэтому $L'$ --- не подрешётка $L$.

\textbf{Диамант $M_3$.} Элементы $a<b,c,d<e$ ($b,c,d$ попарно несравнимы). Проверка недистрибутивности: $b\cdot(c+d)=b\cdot e=b$, но $b\cdot c+b\cdot d=a+a=a\ne b$.

\textbf{Пентагон $N_5$.} Элементы $a<b<c<e$ и $a<d<e$ ($d$ несравним с $b,c$). $c\cdot(b+d)=c\cdot e=c$, но $c\cdot b+c\cdot d=b+a=b\ne c$.

\begin{center}
\begin{tikzpicture}[every node/.style={circle,draw,inner sep=2pt,font=\small},scale=1]
\node (a) at (0,0) {$a$}; \node (b) at (-1,1) {$b$}; \node (c) at (0,1) {$c$}; \node (d) at (1,1) {$d$}; \node (e) at (0,2) {$e$};
\draw (a)--(b)--(e) (a)--(c)--(e) (a)--(d)--(e);
\node[draw=none] at (0,-0.7) {диамант $M_3$};
\begin{scope}[xshift=5cm]
\node (a) at (0,0) {$a$}; \node (b) at (-0.8,0.7) {$b$}; \node (c) at (-0.8,1.5) {$c$}; \node (d) at (0.8,1) {$d$}; \node (e) at (0,2.2) {$e$};
\draw (a)--(b)--(c)--(e) (a)--(d)--(e);
\node[draw=none] at (0,-0.7) {пентагон $N_5$};
\end{scope}
\end{tikzpicture}
\end{center}

\begin{thm}[Биркгоф]
Решётка дистрибутивна тогда и только тогда, когда не имеет подрешёток, изоморфных $M_3$ или $N_5$.
\end{thm}
\textbf{Идея.} ($\Rightarrow$) Подрешётка дистрибутивной дистрибутивна, а $M_3,N_5$ --- нет. ($\Leftarrow$) (в курсе --- без доказательства) если закон нарушен, из нарушающих элементов строится одна из двух подрешёток.

\textbf{Отсутствие псевдодополнений.} В диамант: $b\to c$ был бы наибольшим элементом $\{x\mid b\cdot x\preceq c\}=\{a,c,d\}$ --- наибольшего нет. В пентагоне: $c\to b$: $\{x\mid c\cdot x\preceq b\}=\{a,b,d\}$ --- наибольшего нет. Тем самым диамант и пентагон не импликативны (что согласуется: импликативная $\Rightarrow$ дистрибутивная).

\textbf{Использование.} В теореме о разрешимости ИИВ (билет~\hyperref[T:27]{27}) строится дистрибутивная подрешётка, порождённая значениями подформул; она остаётся алгеброй Гейтинга, хотя псевдодополнения вообще не обязаны сохраняться при переходе к подрешётке.

\begin{warn}
\begin{itemize}
\item «Подрешётка» $\ne$ «подмножество, являющееся решёткой в индуцированном порядке».
\item Диамант $M_3$ и пентагон $N_5$ --- простейшие недистрибутивные решётки (диамант, например, --- решётка из нуля, трёх «прямых» на плоскости над конечным полем и всей плоскости).
\end{itemize}
\end{warn}

% ----------------------------------------------------------------------
\ticket{4(ж)}{Псевдодополнение. Булевы и псевдобулевы алгебры}
\begin{short}
$a\to b=\mathrm{наиб}\{x\mid a\cdot x\preceq b\}$. Импликативная решётка --- в ней существуют все $a\to b$ (имеет 1). Псевдобулева алгебра (алгебра Гейтинга) --- импликативная решётка с $0$, $\sim a=a\to0$. Булева --- ещё и $a+\sim a=1$.
\end{short}

\begin{dfn}
\emph{Псевдодополнение} $a\to b$ --- наибольший элемент множества $\{x\mid a\cdot x\preceq b\}$ (если он существует). \emph{Импликативная решётка} --- решётка, где $a\to b$ существует для любых $a,b$. $0$ --- наименьший элемент решётки, $1$ --- наибольший.
\end{dfn}
\begin{lmm}
В импликативной решётке есть $1$: $a\to a=\mathrm{наиб}\{c\mid a\cdot c\preceq a\}=\mathrm{наиб}\,X=1$.
\end{lmm}
\begin{dfn}
\emph{Псевдобулева алгебра} (алгебра Гейтинга) --- импликативная решётка с $0$. В ней $\sim a:=a\to0$. \emph{Булева алгебра} --- псевдобулева алгебра, в которой $a+\sim a=1$ при всех $a$.
\end{dfn}

\textbf{Свойства импликации} (ДЗ): $a\preceq b\iff a\to b=1$; $b\preceq a\to b$; $a\cdot(a\to b)\preceq b$; $a\preceq b\to c\iff a\cdot b\preceq c$ (присоединение); из $a\preceq b$ следует $b\to c\preceq a\to c$ и $c\to a\preceq c\to b$; $a\to c\preceq(b\to c)\to(a+b\to c)$; $a\to b\preceq(a\to(b\to c))\to(a\to c)$; $a\to(b\to a\cdot b)=1$.

\textbf{Примеры.}
\begin{itemize}
\item $\langle\mathcal{P}(X),\subseteq\rangle$ --- булева алгебра: $a\to b=(X\setminus a)\cup b$, $\sim a=X\setminus a$, $a+\sim a=X$.
\item $\langle\Omega,\subseteq\rangle$ --- псевдобулева алгебра: $a\to b=((X\setminus a)\cup b)^\circ$ (наибольшее открытое, не содержащее точек $a\setminus b$), $\sim a=(X\setminus a)^\circ$. Не булева в общем случае: в $\bbR$ при $a=(0,1)$ имеем $a+\sim a=\bbR\setminus\{0,1\}\ne\bbR$.
\item $\{\text{Л},\text{И}\}$ с $\text{Л}\le\text{И}$ --- булева алгебра в смысле решёток: $a\cdot b=\min$, $a+b=\max$, $a\to b=\neg a\vee b$, $0=\text{Л}$, $1=\text{И}$, $\sim a=\neg a$, $a+\sim a=\text{И}$.
\item Трёхэлементная цепь $0<h<1$: $h\to0=0$, значит $\sim h=0$, $h+\sim h=h\ne1$ --- псевдобулева, не булева (соответствует трёхзначной модели с «неопределённым» значением).
\end{itemize}

\textbf{Согласованность.} Любая псевдобулева алгебра, где $\sem{\cdot}$ подчиняется таблице согласованности, --- корректная модель ИИВ; любая булева --- корректная модель КИВ.

\begin{warn}
\begin{itemize}
\item Псевдобулева $\ne$ булева: аббревиатура <<псевдо>> означает отсутствие закона исключённого третьего.
\item $\sim a$ --- не дополнение, а лишь «псевдодополнение»: $a\cdot\sim a=0$, но $a+\sim a$ может быть $<1$.
\end{itemize}
\end{warn}

% ----------------------------------------------------------------------
\ticket{4(з)}{Алгебра Линденбаума}
\begin{short}
Фактор-алгебра $L/{\approx}$ высказываний по $\alpha\approx\beta\iff\alpha\vd\beta$ и $\beta\vd\alpha$; порядок $[\alpha]\preceq[\beta]\iff\alpha\vd\beta$. Это псевдобулева алгебра, согласованная с ИИВ.
\end{short}

\begin{dfn}
На множестве $L$ всех высказываний определим предпорядок $\alpha\preceq\beta:=\alpha\vd\beta$ (в ИИВ). Положим $\alpha\approx\beta$, если $\alpha\preceq\beta$ и $\beta\preceq\alpha$. \emph{Алгеброй Линденбаума} называется $\mathcal{L}=L/{\approx}$ с порядком на классах $[\alpha]\preceq[\beta]\iff\alpha\preceq\beta$ (корректно: ДЗ).
\end{dfn}
Предпорядок --- рефлексивное и транзитивное отношение ($\alpha\vd\alpha$, и из $\alpha\vd\beta$, $\beta\vd\gamma$ следует $\alpha\vd\gamma$); фактор по $\approx$ делает его частичным порядком.

\begin{thm}
$\mathcal{L}$ --- псевдобулева алгебра, причём оценка $\sem{\alpha}=[\alpha]_{\mathcal{L}}$ согласована с ИИВ.
\end{thm}
\begin{proof}[Схема]
Надо показать: $\preceq$ --- порядок на $\mathcal{L}$ и
\[
[\alpha\vee\beta]=[\alpha]+[\beta],\ \ [\alpha\amp\beta]=[\alpha]\cdot[\beta],\ \ [\alpha\to\beta]=[\alpha]\to[\beta],\ \ [A\amp\neg A]=0,\ \ [A\to A]=1,\ \ [\neg\alpha]=\ \sim[\alpha].
\]
\begin{itemize}
\item \emph{Сумма и произведение}: из аксиом 4--8: $\alpha\amp\beta\vd\alpha$, $\alpha\amp\beta\vd\beta$ (нижняя грань); если $\sigma\vd\alpha$ и $\sigma\vd\beta$, то $\sigma\vd\alpha\amp\beta$ (сх.\ 3 и дедукция) --- наибольшая. Аналогично для $\vee$.
\item \emph{Импликация}: $[\alpha]\to[\beta]=\mathrm{наиб}\{[\sigma]\mid[\alpha\amp\sigma]\preceq[\beta]\}$. Два включения: $\alpha\amp(\alpha\to\beta)\vd\beta$ (карринг и M.P.); если $\alpha\amp\sigma\vd\beta$, то $\sigma\vd\alpha\to\beta$ (карринг и теорема о дедукции).
\item \emph{Ноль}: $0=[A\amp\neg A]$, так как $A\amp\neg A\vd\sigma$ при любом $\sigma$ (принцип взрыва).
\item \emph{Отрицание}: $[\neg\alpha]=[\alpha\to A\amp\neg A]=[\alpha]\to0=\ \sim[\alpha]$; $[A\to A]=[A]\to[A]=1$.\qedhere
\end{itemize}
\end{proof}

\textbf{Роль.} Алгебра Линденбаума «переворачивает» синтаксис в семантику: доказуемые формулы образуют класс $1_{\mathcal{L}}$ (так как $\vd\alpha\iff[\alpha]=[A\to A]$). Именно её используют для полноты (билет~\hyperref[T:23]{23}), для дизъюнктивности (билет~\hyperref[T:31]{31}) и для разрешимости.
\emph{Предостережение:} в КИВ алгебра Линденбаума --- булева.

\begin{warn}
\begin{itemize}
\item Элементы $\mathcal{L}$ --- классы формул, а не формулы; в $\mathcal{L}$ бесконечно много элементов.
\item Порядок определён через выводимость в \emph{интуиционистском} исчислении; в классическом получилась бы булева алгебра.
\end{itemize}
\end{warn}

% ----------------------------------------------------------------------
\ticket{4(и)}{Полнота ИИВ в псевдобулевых алгебрах (формулировка, идея доказательства)}
\begin{short}
Если $\alpha$ истинна во всех псевдобулевых алгебрах, то $\vd_{\text{и}}\alpha$. Доказательство: возьмём в качестве модели алгебру Линденбаума; $\sem{\alpha}=1_{\mathcal{L}}$ означает $\vd\alpha$.
\end{short}

\begin{thm}[корректность]
Любая псевдобулева алгебра, являющаяся согласованной оценкой ИИВ, --- корректная модель: $\vd\alpha\Rightarrow\sem{\alpha}=1$. (Индукция по выводу; свойства импликативных решёток из ДЗ дают общезначимость всех 10 схем в варианте 10и и сохранение M.P.)
\end{thm}
\begin{thm}[полнота]
Интуиционистское исчисление высказываний полно в псевдобулевых алгебрах: если $\models\alpha$ во всех псевдобулевых алгебрах, то $\vd\alpha$.
\end{thm}
\begin{proof}
Возьмём в качестве модели алгебру Линденбаума: $\sem{\alpha}=[\alpha]_{\mathcal{L}}$ (оценка согласована, билет~\hyperref[T:22]{22}). Пусть $\alpha$ истинна во всех псевдобулевых алгебрах, в том числе в $\mathcal{L}$: $\sem{\alpha}=1_{\mathcal{L}}$, т.\,е. $[\alpha]_{\mathcal{L}}=[A\to A]_{\mathcal{L}}$. Значит $A\to A\approx\alpha$, в частности $A\to A\vd\alpha$. Но $\vd A\to A$, поэтому $\vd\alpha$.
\end{proof}
Идея: «идейно мы просто перекидываем всё в алгебру Линденбаума и пользуемся её свойствами».

\textbf{Следствия и связь.}
\begin{itemize}
\item Топологические модели: $\models\alpha$ в топологических моделях (при всех $\langle X,\Omega\rangle$ значение равно $X$) $\iff\vd_{\text{и}}\alpha$.
\item Аналогично в КИВ: булева алгебра $\Leftrightarrow$ полнота в булевых алгебрах.
\item ИИВ корректно и полно в псевдобулевых алгебрах, но \emph{не} в конечных табличных моделях (билет~\hyperref[T:26]{26}) --- однако для каждой недоказуемой формулы существует конечная алгебра Гейтинга (билет~\hyperref[T:27]{27}).
\end{itemize}

\textbf{Сводка общих результатов.}
\begin{center}
\begin{tabular}{l|l|l}
&КИВ&ИИВ (и алгебры Гейтинга)\\\hline
корректность&да (лекция 1)&да (ДЗ 3)\\
непротиворечивость&да&да\\
полнота&да (лекция 2)&да (лекция 3)\\
разрешимость&да (таблицы)&да (лекция 4)
\end{tabular}
\end{center}

\begin{warn}
\begin{itemize}
\item Полнота формулируется «во всех псевдобулевых алгебрах» (на лекции) или «во всех топологических моделях»; о достаточности одного пространства (например $\bbR$ с евклидовой топологией) на лекции не говорится.
\item Шаг «$\alpha\approx A\to A\Rightarrow\vd\alpha$» нужен именно потому, что в $\mathcal{L}$ единица --- класс всех доказуемых формул.
\end{itemize}
\end{warn}

% ----------------------------------------------------------------------
\ticket{4(к)}{Модели Крипке. Вынужденность}
\begin{short}
$\langle\mathcal{W},\preceq,\Vdash\rangle$: миры, частичный порядок, вынужденность переменных монотонна. $W\Vdash\alpha\to\beta$ --- во всех $W_1\succeq W$ из $W_1\Vdash\alpha$ следует $W_1\Vdash\beta$; $W\Vdash\neg\alpha$ --- нигде выше не вынуждается $\alpha$.
\end{short}

\textbf{Идея.} Миры --- возможные состояния знания; $W_1\preceq W_2$ --- состояние $W_2$ достижимо из $W_1$ (мы узнали больше). Научное знание: факты, которых мы пока не знаем, когда-нибудь будут известны, возможно как «да», возможно как «нет» (напр., $P=NP$?). В рассуждениях стоит учитывать оба варианта.

\begin{dfn}
\emph{Модель Крипке} --- тройка $\langle\mathcal{W},(\preceq),(\Vdash)\rangle$:
\begin{itemize}
\item $\mathcal{W}$ --- множество миров, $\preceq$ --- нестрогий частичный порядок на $\mathcal{W}$;
\item $\Vdash\ \subseteq\mathcal{W}\times\mathcal{P}$ --- отношение \emph{вынужденности} миров и переменных, монотонное: если $W_i\preceq W_j$ и $W_i\Vdash X$, то $W_j\Vdash X$.
\end{itemize}
Вынужденность доопределяется на формулы:
\begin{itemize}
\item $W\Vdash\alpha\amp\beta$, если $W\Vdash\alpha$ и $W\Vdash\beta$;
\item $W\Vdash\alpha\vee\beta$, если $W\Vdash\alpha$ или $W\Vdash\beta$;
\item $W\Vdash\alpha\to\beta$, если для всех $W_1\succeq W$ из $W_1\Vdash\alpha$ следует $W_1\Vdash\beta$;
\item $W\Vdash\neg\alpha$, если для всех $W_1\succeq W$ неверно $W_1\Vdash\alpha$.
\end{itemize}
Пишем $\Vdash\alpha$, если $W\Vdash\alpha$ для всех $W\in\mathcal{W}$; $\models_{\kappa}\alpha$, если $\Vdash\alpha$ во всех моделях Крипке.
\end{dfn}

\begin{lmm}[монотонность]
Если $W_1\Vdash\alpha$ и $W_1\preceq W_2$, то $W_2\Vdash\alpha$.
\end{lmm}
\begin{proof}
Индукция по $\alpha$. Переменная --- по определению; $\amp,\vee$ --- очевидно; для $\to$ и $\neg$: все миры, достижимые из $W_2$, достижимы из $W_1$ (транзитивность $\preceq$).
\end{proof}

\textbf{Пример: ЗИТ.} Модель: $W_1\preceq W_2$, $W_1\preceq W_3$, $W_2\Vdash A$, $W_3\not\Vdash A$.
\begin{center}
\begin{tikzpicture}[node distance=1.6cm,every node/.style={font=\small}]
\node (w1) at (0,0) {$W_1$};
\node (w2) at (2.4,0.9) {$W_2\ (\Vdash A)$};
\node (w3) at (2.4,-0.9) {$W_3$};
\draw[-{Stealth}] (w1)--(w2); \draw[-{Stealth}] (w1)--(w3);
\end{tikzpicture}
\end{center}
$W_3\Vdash\neg A$ (выше $W_3$ нет миров с $A$), но $W_1\not\Vdash A$ и $W_1\not\Vdash\neg A$ (так как $W_1\preceq W_2\Vdash A$). Значит $W_1\not\Vdash A\vee\neg A$, откуда $\not\models_\kappa A\vee\neg A$.

\textbf{Корректность.} Любая модель Крипке --- корректная модель ИИВ (билет~\hyperref[T:25]{25}); верна и теорема о полноте относительно конечных деревьев (опровергаемость конечной древовидной моделью). Миры с одним узлом дают классические модели.

\textbf{Другие опровержения (ДЗ 4).} В моделях Крипке опровергаются закон Пирса $((A\to B)\to A)\to A$, $(A\to B)\to\neg A\vee B$, $(A\to(B\vee\neg B))\vee(\neg A\to(B\vee\neg B))$.

\begin{warn}
\begin{itemize}
\item Дизъюнкция истинна в мире, если истинен \emph{один из членов в нём же}, а не «где-то выше» --- иначе исчезло бы дизъюнктивное свойство.
\item Для $\neg\alpha$ нужно смотреть на все миры $\succeq W$, включая сам $W$.
\item В исходной формулировке монотонность требуется \emph{только} для переменных; для составных формул она --- лемма.
\end{itemize}
\end{warn}

% ----------------------------------------------------------------------
\ticket{4(л)}{Сведение моделей Крипке к псевдобулевым алгебрам}
\begin{short}
Для древовидной (вообще любой) модели Крипке множества $V(\alpha)=\{w\mid w\Vdash\alpha\}$ открыты в топологии верхних множеств; открытые множества образуют алгебру Гейтинга, и $\sem{\alpha}=V(\alpha)$ --- её согласованная оценка.
\end{short}

\begin{thm}
Любая модель Крипке --- корректная модель ИИВ.
\end{thm}
\begin{proof}
Пусть порядок древовидный (для произвольного --- аналогично с топологией Александрова). Топология на деревьях: множество открыто, если вместе с вершиной содержит всех её потомков, т.\,е. $S$ открыто $\iff$ $S$ замкнуто вверх по $\preceq$.
\begin{enumerate}
\item $V(\alpha):=\{w\in\mathcal{W}\mid w\Vdash\alpha\}$ открыто --- это лемма о монотонности.
\item Положим $V=\{S\subseteq\mathcal{W}\mid S\text{ открыто}\}$ --- псевдобулева алгебра (билет~\hyperref[T:21]{21}) относительно $\subseteq$ и $\sem{\alpha}:=V(\alpha)$.
\item Согласованность: $V(\alpha\amp\beta)=V(\alpha)\cap V(\beta)$, $V(\alpha\vee\beta)=V(\alpha)\cup V(\beta)$.
  $V(\alpha\to\beta)=\{W\mid\forall W_1\succeq W:\ W_1\in V(\alpha)\Rightarrow W_1\in V(\beta)\}$ --- это в точности наибольшее открытое (верхнее) множество, содержащееся в $(\mathcal{W}\setminus V(\alpha))\cup V(\beta)$, т.\,е. $\big(\cmpl V(\alpha)\cup V(\beta)\big)^\circ$. Аналогично $V(\neg\alpha)=(\cmpl V(\alpha))^\circ$.
\end{enumerate}
Значит оценка $V(\cdot)$ --- согласованная оценка алгеброй Гейтинга, а поэтому (билет~\hyperref[T:23]{23}) корректна: $\vd\alpha\Rightarrow\sem{\alpha}=\mathcal{W}$, т.\,е. $\Vdash\alpha$.
\end{proof}

\textbf{Смысл.} Внутренность и дополнение топологической модели --- ровно «для всех будущих миров» модели Крипке. Модели Крипке --- частный случай топологических: пространство $\langle\mathcal{W},\{\text{верхние множества}\}\rangle$.

\textbf{Обратное.} Если формула опровергается в некоторой алгебре Гейтинга, то она опровергается и в некоторой (конечной, древовидной) модели Крипке --- полнота ИИВ относительно моделей Крипке (билеты по лекции 4--5; в коллоквиуме требуется только сведение).

\begin{warn}
\begin{itemize}
\item Топология на деревьях здесь «верхняя»: открыты множества, замкнутые вверх (содержат потомков).
\item В общем случае (не дерево) нужно говорить «верхние множества предпорядка» (топология Александрова); доказательство то же.
\end{itemize}
\end{warn}

% ----------------------------------------------------------------------
\ticket{4(м)}{Нетабличность ИИВ (формулировка теоремы)}
\begin{short}
Не существует полной конечной табличной модели для ИИВ (конечного набора значений с функциями связок, доказывающей ровно то, что доказуемо). Доказательство через формулы $\alpha_n=\bigvee_{p<q}(A_p\to A_q)$.
\end{short}

\begin{dfn}[табличная модель]
Заданы множество $V$, значение $T\in V$ («истина»), функции $f_{\amp},f_\vee,f_\to:V\times V\to V$, $f_\neg:V\to V$ и оценка $f_P:\mathcal{P}\to V$. Оценка $\sem{X}=f_P(X)$, $\sem{\alpha\star\beta}=f_\star(\sem{\alpha},\sem{\beta})$, $\sem{\neg\alpha}=f_\neg(\sem{\alpha})$ --- \emph{табличная}. Если $\vd\alpha$ влечёт $\sem{\alpha}=T$ при всех $f_P$, то $\mathcal{M}=\langle V,T,f_{\amp},f_\vee,f_\to,f_\neg\rangle$ --- \emph{табличная модель}; она \emph{конечна}, если $V$ конечно; \emph{полна}, если верно и обратное ($\sem{\alpha}=T$ при всех $f_P$ влечёт $\vd\alpha$).
\end{dfn}
\begin{thm}
Не существует полной конечной табличной модели для интуиционистского исчисления высказываний.
\end{thm}

\textbf{Схема доказательства.} Пусть $V=\{v_1,\dots,v_n\}$ --- конечная полная модель. Рассмотрим
\[
\alpha_n=\bigvee_{1\le p<q\le n+1}(A_p\to A_q).
\]
\begin{enumerate}
\item В модели $\alpha_n$ \emph{общезначима}: по принципу Дирихле для любой оценки $A_1,\dots,A_{n+1}\to V$ найдутся $p<q$ с $\sem{A_p}=\sem{A_q}=v$. Тогда $\sem{A_p\to A_q}=f_\to(v,v)=T$ (так как $\vd X\to X$ и модель корректна). Дизъюнкция $\sigma\vee(X\to X)\vee\tau$ принимает значение $T$ ($\vd\sigma\vee(X\to X)\vee\tau$). Значит $\sem{\alpha_n}=T$.
\item Но $\nvd\alpha_n$: в модели Крипке с корнем $W_R$ (в нём ничего не вынуждается) и $n$ мирами-листьями $W_1,\dots,W_n$, где $W_i\Vdash A_j\iff j\le i$ (так что $W_i\Vdash A_i$, но $W_i\not\Vdash A_q$ при $q>i$; $A_{n+1}$ не вынуждается нигде), имеем $W_p\not\Vdash A_p\to A_q$ при $p<q$ (сам $W_p\succeq W_p$ вынуждает $A_p$, но не $A_q$), значит $W_R\not\Vdash A_p\to A_q$ при всех $p<q$ и $W_R\not\Vdash\alpha_n$.
\end{enumerate}
Итак, $\alpha_n$ общезначима в конечной модели, но недоказуема --- модель неполна. (Для корректной формализации используются факты ДЗ~4: $\sem{(\alpha\to\alpha)\vee\beta}=1$, $\sem{\alpha\vee\beta}=\sem{\beta\vee\alpha}$ и $\sem{(\alpha\vee\beta)\vee\gamma}=\sem{\alpha\vee(\beta\vee\gamma)}$ --- они нужны, чтобы вынести из дизъюнкции члены с совпадающими значениями.)

\begin{center}
\begin{tikzpicture}[every node/.style={font=\small}]
\node (r) at (0,0) {$W_R$};
\node (w1) at (3,1.3) {$W_1\ (\Vdash A_1)$};
\node (w2) at (3,0.4) {$W_2\ (\Vdash A_2)$};
\node (w3) at (3,-0.4) {$\cdots$};
\node (wn) at (3,-1.3) {$W_n\ (\Vdash A_n)$};
\draw[-{Stealth}] (r)--(w1); \draw[-{Stealth}] (r)--(w2); \draw[-{Stealth}] (r)--(wn);
\end{tikzpicture}
\end{center}

\textbf{Смысл.} Конечными таблицами интуиционистскую логику не охарактеризовать (в отличие от классической: достаточно двух значений). Однако ИИВ разрешимо (билет~\hyperref[T:27]{27}): нужно \emph{семейство} конечных моделей, растущее с размером формулы.

\begin{warn}
\begin{itemize}
\item На коллоквиуме требуется только \emph{формулировка}; блок «Схема» --- на случай уточняющих вопросов.
\item «Нетабличность» не означает неразрешимость.
\item В $\Omega(\bbR)$ формулу $\alpha_n$ можно опровергнуть, выбрав $n+1$ различных открытых множеств с подходящим порядком включения (ДЗ 4).
\end{itemize}
\end{warn}

% ----------------------------------------------------------------------
\ticket{4(н)}{Разрешимость ИИВ (формулировка теорем, конструкция-аналог ДНФ на решётках)}
\begin{short}
Если $\nvd\alpha$, то есть алгебра Гейтинга $\mathcal{G}\not\models\alpha$ размера $\le2^{2^{|\alpha|+2}}$. Достаточно перебрать все такие алгебры. Ключ: дистрибутивная подрешётка, порождённая $n$ элементами, содержит $\le2^{2^n}$ элементов (аналог ДНФ).
\end{short}

\begin{dfn}
Язык $\mathcal{L}\subseteq A^*$ \emph{разрешим}, если есть алгоритм, завершающийся на любом $a\in A^*$ и выдающий <<истина>> при $a\in\mathcal{L}$ и <<ложь>> иначе. Теория разрешима, если разрешим язык доказуемых формул. Пример неразрешимого: язык программ машин Тьюринга, останавливающихся на любом входе.
\end{dfn}

\begin{lmm}
Существует дистрибутивная подрешётка $\mathcal{L}'$, содержащая $a_1,\dots,a_n$, что $|L'|\le2^{2^n}$.
\end{lmm}
\begin{proof}
Пусть $L'=\{\varphi(a_1,\dots,a_n)\mid\varphi$ составлено из $(+)$ и $(\cdot)\}$ (замыкание порождающих относительно $+$ и $\cdot$). Для $p,q\in L'$ имеем $p\star q\in L'$ (так как $\varphi_p(\vec a)\star\varphi_q(\vec a)=\psi(\vec a)$), и ясно, что $+$ и $\cdot$ в $L'$ совпадают с $+$ и $\cdot$ в $\mathcal{L}$; значит $\mathcal{L}'$ --- подрешётка, и она дистрибутивна. Из дистрибутивности каждый элемент приводится к «ДНФ»
\[
\varphi(a_1,\dots,a_n)=\sum_{\text{Кн}\in\text{ДНФ}(\varphi)}\ \prod_{i\in\text{Кн}}a_i
\]
(сумма произведений подмножеств порождающих). Конъюнктов (подмножеств $\{a_i\}$) не больше $2^n$; различных сумм конъюнктов (подмножеств множества конъюнктов) не больше $2^{2^n}$.
\end{proof}

\begin{thm}
Если $\nvd\alpha$ в ИИВ, то существует алгебра Гейтинга $\mathcal{G}$, что $\mathcal{G}\not\models\alpha$, причём $|\mathcal{G}|\le2^{2^{|\alpha|+2}}$.
\end{thm}
\begin{proof}
Если $\nvd\alpha$, то по полноте найдётся алгебра Гейтинга $\mathcal{H}$ с $\mathcal{H}\not\models\alpha$. Пусть $\varphi_1,\dots,\varphi_n$ --- все подформулы $\alpha$ ($n\le|\alpha|$). Пусть $\mathcal{G}$ --- дистрибутивная подрешётка $\mathcal{H}$, порождённая $\sem{\varphi_1},\dots,\sem{\varphi_n}$, а также $0$ и $1$ (по лемме $|G|\le2^{2^{n+2}}$). $\mathcal{G}$ --- конечная дистрибутивная решётка, значит импликативна, т.\,е. алгебра Гейтинга. Значения подформул остаются теми же (по индукции: псевдодополнение $a\to_{\mathcal H}b$ подформул лежит в $G$ и остаётся наибольшим в $G$), значит $\mathcal{G}\not\models\alpha$.
\end{proof}
\begin{thm}
ИИВ разрешимо.
\end{thm}
\begin{proof}
По формуле $\alpha$ перечисляем все алгебры Гейтинга $\mathcal{G}$ размера $\le2^{2^{|\alpha|+2}}$ (конечно много) и все оценки переменных. Если $\alpha$ истинна во всех, то $\vd\alpha$ (иначе нашлась бы опровергающая малая алгебра); если найдена опровергающая --- $\nvd\alpha$ (по корректности).
\end{proof}

\textbf{Замечания.} Сложность алгоритма огромна; в источниках (конспект Чепелина) использован запас $2^{2^{|\alpha|+37}}$, но идея та же. Разрешимость обеспечивает «конечная модельная» (finite model) свойство ИИВ.

\begin{warn}
\begin{itemize}
\item Не путать с нетабличностью: ни одна \emph{единственная} конечная таблица не достаточна, но семейство таблиц ограниченного размера (зависящего от $|\alpha|$) --- достаточно.
\item На коллоквиуме требуются формулировки теорем и пояснение конструкции ДНФ: «любой элемент подрешётки --- сумма произведений порождающих».
\end{itemize}
\end{warn}

% ======================================================================
\group{5. Изоморфизм Карри--Ховарда}{Билет 28.\par\hyperref[T:28]{28}}

\ticket{5}{Изоморфизм Карри--Ховарда: интерпретация связок. Какие конструкции имеют тип «ложь»?}
\begin{short}
Доказательство = программа, высказывание = тип, импликация = функция, $\amp$ = пара, $\vee$ = тип-сумма (variant), $\bot$ = необитаемый тип. Значения типа «ложь» (с оговорками) --- выражения, не возвращающие управления: исключения, бесконечные циклы, \texttt{abort}.
\end{short}

\textbf{Идея.} BHK-интерпретация говорит о <<построениях>>. Пусть построение --- значение типа, которое может быть вычислено. Тогда
\begin{itemize}
\item $\alpha\amp\beta$ --- тип упорядоченной пары значений $\alpha\times\beta$;
\item $\alpha\vee\beta$ --- алгебраический тип (тип-сумма) --- значение типа $\alpha$ \emph{или} типа $\beta$, и мы знаем, какого именно (\texttt{std::variant} в C++, \texttt{Either a b} в Haskell);
\item $\alpha\to\beta$ --- тип функций из $\alpha$ в $\beta$;
\item $\bot$ --- тип без значений (необитаемый).
\end{itemize}
Доказательство --- значение, \emph{обитающее} в этом типе: $\vd3:\texttt{int}$, $\vd\texttt{fun x -> x+1}:\texttt{int}\to\texttt{int}$.

\begin{center}
\begin{tabular}{l|l}
Программа ($\lambda$-выражение)&Исчисление высказываний\\\hline
Выражение&Доказательство\\
Тип выражения&Высказывание\\
Тип функции&Импликация\\
Упорядоченная пара&Конъюнкция\\
Алгебраический тип (сумма)&Дизъюнкция\\
Необитаемый тип&Ложь
\end{tabular}
\end{center}
Обычно изоморфизм формулируется для $\lambda$-исчисления (формальный алгоритм, абстрагирующий синтаксис языков): типизированное $\lambda$-исчисление $\leftrightarrow$ натуральный вывод ИИВ.
Правила вывода соответствуют конструкциям: введение $\to$ --- $\lambda$-абстракция, удаление $\to$ --- применение функции; введение $\amp$ --- конструктор пары, удаление --- проекции; введение $\vee$ --- конструкторы \texttt{Left}/\texttt{Right}; удаление $\vee$ --- \texttt{match}/\texttt{case}.

\textbf{Тип-сумма (дизъюнкция).} Отмеченное объединение $A\uplus B=\{\langle x,0\rangle\mid x\in A\}\cup\{\langle y,1\rangle\mid y\in B\}$. В OCaml:
\begin{verbatim}
type 'a option = Some of 'a | None      (* 'a option = 'a \/ ()  *)
let csqrt x = if x >= 0. then Some (sqrt x) else None
\end{verbatim}

\textbf{Ложь.} Перепишем:
\begin{verbatim}
let csqrt x =
   if x >= 0. then sqrt x
              else failwith "Cannot compute square root"
\end{verbatim}
Ветка \texttt{then} имеет тип \texttt{float}; \texttt{failwith s} --- выражение, не возвращающее результата, и поэтому его тип можно считать любым: $\text{failwith s}:\bot\ \vd\ \text{failwith s}:\texttt{float}$. Это и есть <<из лжи следует всё, что угодно>> (принцип взрыва).

\textbf{Какие конструкции имеют тип «ложь»?} Выражения, которые \emph{никогда не возвращают значение}:
\begin{itemize}
\item исключения (\texttt{throw}, \texttt{raise}, \texttt{failwith});
\item бесконечные циклы и бесконечная рекурсия (\texttt{let rec f x = f x});
\item аварийное завершение (\texttt{abort}, \texttt{exit}, \texttt{std::terminate});
\item \texttt{undefined}, \texttt{error "..."} в Haskell;
\item в языках с типом \texttt{Void}/\texttt{Empty}: функция $\bot\to A$ (\texttt{absurd}) есть всегда, значения $\bot$ нет.
\end{itemize}
Их наличие делает типовую систему \emph{противоречивой} (в ней «доказуемо» $\bot$), поэтому в Coq/Agda (логически корректных) такие конструкции запрещены: все функции тотальны.

\textbf{Отрицание} $\neg A=A\to\bot$: функция, из любого значения $A$ дающая невозможное. Закон $A\vee\neg A$ означал бы программу, по любому типу $A$ \emph{решающую}, обитаем ли он --- такой нет.

\textbf{Классический вариант.} Закон Пирса (\texttt{call/cc}) и $\neg\neg A\to A$ соответствуют операторам управления (continuations).

\begin{warn}
\begin{itemize}
\item $\bot$ --- не тип, содержащий значение «ошибка»: именно \emph{необитаемый} тип. Выражение, зацикливающееся, «имеет» любой тип лишь потому, что не возвращается.
\item Изоморфизм не используют для доказательства «из программ»: нужно, чтобы программа не использовала исключения и нетерминирующие конструкции (ДЗ 3).
\item Алгебраический тип $\alpha\vee\beta$ содержит \emph{метку} --- вот почему $\alpha\vee\neg\alpha$ нельзя построить в общем случае.
\end{itemize}
\end{warn}

% ======================================================================
\group{6. Дизъюнктивность интуиционистского исчисления высказываний}{Билеты 29--31.\par\hyperref[T:29]{29}\ \hyperref[T:30]{30}\ \hyperref[T:31]{31}}

\ticket{6(а)}{Гёделева алгебра. Операция \texorpdfstring{$\Gamma(\mathcal{A})$}{Γ(A)}}
\begin{short}
Решётка гёделева, если $a+b=1\Rightarrow a=1$ или $b=1$. Операция $\Gamma(\mathcal{A})$ добавляет к алгебре Гейтинга новый элемент $\omega$ непосредственно под единицей. $\Gamma(\mathcal{A})$ --- гёделева алгебра.
\end{short}

\begin{dfn}
Решётка \emph{гёделева}, если $a+b=1$ влечёт $a=1$ или $b=1$. (Дизъюнктивность на уровне алгебры.)
\end{dfn}
\begin{dfn}[гёделевизация]
Для алгебры Гейтинга $\mathcal{A}=\langle A,\preceq\rangle$ определим $\Gamma(\mathcal{A})=\langle A\cup\{\omega\},\preceq_{\Gamma(\mathcal{A})}\rangle$, где $\omega\notin A$ --- новый элемент, а $\preceq_{\Gamma(\mathcal{A})}$ --- минимальное отношение порядка, удовлетворяющее:
\begin{itemize}
\item $a\preceq_{\Gamma(\mathcal{A})}b$, если $a\preceq_{\mathcal{A}}b$ и $a,b\notin\{\omega,1\}$;
\item $a\preceq_{\Gamma(\mathcal{A})}\omega$, если $a\ne1$;
\item $\omega\preceq_{\Gamma(\mathcal{A})}1$.
\end{itemize}
\end{dfn}
Картинка: старая алгебра без единицы (серый блок), над ней новый элемент $\omega$, над ним старая единица $1$.
\begin{center}
\begin{tikzpicture}[every node/.style={font=\small}]
\fill[gray!25] (0,0) circle (1.1);
\node at (0,0) {$A\setminus\{1\}$};
\node[circle,fill,inner sep=1.5pt,label=right:$\omega$] (o) at (0,1.5) {};
\node[circle,fill,inner sep=1.5pt,label=right:$1$] (t) at (0,2.5) {};
\draw[-{Stealth}] (0,1.1)--(o); \draw[-{Stealth}] (o)--(t);
\end{tikzpicture}
\end{center}

\begin{thm}
$\Gamma(\mathcal{A})$ --- гёделева алгебра (и алгебра Гейтинга).
\end{thm}
\begin{proof}
\emph{Гёделева}: пусть $a+b=1$ и $a\ne1$, $b\ne1$. Тогда $a\preceq\omega$ и $b\preceq\omega$, а значит $a+b\preceq\omega\prec1$ --- противоречие. Следовательно хотя бы одно из $a,b$ равно $1$.

\emph{Алгебра Гейтинга}: порядок --- решётка: для $a,b\in A\setminus\{1\}$ супремум в $A$ остаётся супремумом, если он $\ne1$; если $a+_{\mathcal{A}}b=1$, то в $\Gamma(\mathcal{A})$ общих верхних граней, кроме $\omega$ и $1$, нет (любая $c\in A\setminus\{1\}$ с $a,b\preceq c$ дала бы $c\succeq a+b=1$), так что $a+b=\omega$; $0$ остаётся нулём. Псевдодополнения: для $a,b\in A\setminus\{1\}$ имеем $a\to b=a\to_{\mathcal{A}}b$ (при $a\preceq b$ это $1$, иначе $\ne1$, и $\omega$ в множество $\{x\mid a\cdot x\preceq b\}$ не входит, так как $a\cdot\omega=a\not\preceq b$); $a\to\omega=\omega\to\omega=1$; $1\to b=b$; $\omega\to b=b$ для $b\in A\setminus\{1\}$; $1\to\omega=\omega$.
\end{proof}

\textbf{Оценка в $\Gamma(\mathcal{L})$.} Напомним, $\mathcal{L}$ --- алгебра Линденбаума. Положим $\sem{X}_{\Gamma(\mathcal{L})}:=\sem{X}_{\mathcal{L}}$ и связки оцениваем естественно: $\sem{\alpha\amp\beta}_{\Gamma(\mathcal{L})}=\sem{\alpha}_{\Gamma(\mathcal{L})}\cdot\sem{\beta}_{\Gamma(\mathcal{L})}$ и т.\,п. Эта оценка является алгеброй Гейтинга, согласованной с ИИВ: $\sem{\alpha\to\alpha}=1$, $\sem{\alpha\amp\neg\alpha}=0$ и т.\,д.

\textbf{Вопрос ДЗ: бывает ли $\sem{\alpha}_{\Gamma(\mathcal{L})}=\omega$?} Нет, ни для какой формулы при этой оценке: $\mathcal{G}(\omega)=1$, поэтому из $\sem{\alpha}_{\Gamma(\mathcal{L})}=\omega$ следовало бы $\sem{\alpha}_{\mathcal{L}}=1$, т.\,е. $\vd\alpha$, а доказуемые формулы получают значение $1$, а не $\omega$ (согласованность). Элемент $\omega$ в образе оценки не встречается; он нужен лишь для того, чтобы алгебра стала гёделевой.

\begin{warn}
\begin{itemize}
\item $\omega$ --- \emph{не} «бесконечность»; это просто имя нового элемента.
\item В $\Gamma(\mathcal{A})$ новая единица отделяет «почти истинное» ($\omega$) от «истинного» ($1$); именно это и даёт дизъюнктивное свойство.
\end{itemize}
\end{warn}

% ----------------------------------------------------------------------
\ticket{6(б)}{Гомоморфизм алгебр}
\begin{short}
Гомоморфизм алгебр Гейтинга сохраняет $+,\cdot,\to,0,1$. Отображение $\mathcal{G}:\Gamma(\mathcal{L})\to\mathcal{L}$ ($\omega\mapsto1$, остальное --- тождественно) --- гомоморфизм, согласованный с оценками.
\end{short}

\begin{dfn}
Пусть $\mathcal{A},\mathcal{B}$ --- алгебры Гейтинга. Отображение $g:\mathcal{A}\to\mathcal{B}$ --- \emph{гомоморфизм}, если $g(a\star b)=g(a)\star g(b)$ для $\star\in\{+,\cdot,\to\}$, $g(0_{\mathcal{A}})=0_{\mathcal{B}}$ и $g(1_{\mathcal{A}})=1_{\mathcal{B}}$.
\end{dfn}
\begin{dfn}
Оценка $\sem{\cdot}_{\mathcal{A}}$ \emph{согласована} с $\sem{\cdot}_{\mathcal{B}}$ и гомоморфизмом $g$, если $g(\mathcal{A})=\mathcal{B}$ (сюръективность) и $g(\sem{\alpha}_{\mathcal{A}})=\sem{\alpha}_{\mathcal{B}}$.
\end{dfn}

\textbf{Гомоморфизм Гёделя.} Определим $\mathcal{G}:\Gamma(\mathcal{L})\to\mathcal{L}$:
\[
\mathcal{G}(a)=\begin{cases}a,&a\ne\omega\\1,&a=\omega.\end{cases}
\]
\begin{lmm}
$\mathcal{G}$ --- гомоморфизм $\Gamma(\mathcal{L})$ в $\mathcal{L}$, причём оценка $\sem{\cdot}_{\Gamma(\mathcal{L})}$ согласована с $\mathcal{G}$ и $\sem{\cdot}_{\mathcal{L}}$.
\end{lmm}
\begin{proof}[Схема]
Проверка по случаям. Сумма: если $a+_{\mathcal{L}}b=1$ при $a,b\ne\omega$, то в $\Gamma(\mathcal{L})$ сумма равна $\omega$ (если оба $\ne1$) или $1$, и $\mathcal{G}$ даёт $1=a+b$; если $a=\omega$, то $\mathcal{G}(\omega+b)=1=1+b$. Произведение: $\omega\cdot a=a$ при $a\ne1$, и $\mathcal{G}(a)=a=1\cdot a$. Импликация: $\omega\to b=b$ и $\mathcal{G}(\omega)\to\mathcal{G}(b)=1\to b=b$; $a\to\omega=1=a\to1$; $1\to\omega=\omega\mapsto1=1\to1$; для $a,b\ne\omega$ значения $a\to b$ совпадают. Нуль и единица сохраняются ($\mathcal{G}(1)=1$, $\mathcal{G}(0)=0$). Оценка согласована индукцией по формуле: переменные --- по определению, связки --- по гомоморфности.
\end{proof}
Таким образом $\mathcal{G}$ <<склеивает>> $\omega$ с единицей.

\begin{warn}
\begin{itemize}
\item Гомоморфизм сохраняет $0$ и $1$ --- это обязательные условия определения.
\item Образ гомоморфизма --- подалгебра; сюръективность нужна для «согласованности», чтобы значение любой формулы в $\mathcal{B}$ имело прообраз.
\end{itemize}
\end{warn}

% ----------------------------------------------------------------------
\ticket{6(в)}{Дизъюнктивность ИИВ (формулировка, идея доказательства)}
\begin{short}
Если $\vd_{\text{и}}\alpha\vee\beta$, то $\vd\alpha$ или $\vd\beta$. Идея: оценить в $\Gamma(\mathcal{L})$ ($=1$ $\Rightarrow$ один из членов $=1$, ибо гёделева), спустить гомоморфизмом $\mathcal{G}$ в $\mathcal{L}$ и применить полноту $\mathcal{L}$.
\end{short}

\begin{dfn}
Исчисление \emph{дизъюнктивно}, если при любых $\alpha,\beta$ из $\vd\alpha\vee\beta$ следует $\vd\alpha$ или $\vd\beta$.
\end{dfn}
\begin{thm}[дизъюнктивность ИИВ]
Если $\vd_{\text{и}}\alpha\vee\beta$, то $\vd_{\text{и}}\alpha$ или $\vd_{\text{и}}\beta$.
\end{thm}
\begin{proof}
Пусть $\vd\alpha\vee\beta$.
\begin{enumerate}
\item Тогда $\sem{\alpha\vee\beta}_{\Gamma(\mathcal{L})}=1$ (оценка согласована с ИИВ, а доказуемые формулы принимают значение $1$).
\item Так как $\sem{\alpha\vee\beta}=\sem{\alpha}+\sem{\beta}=1$ и $\Gamma(\mathcal{L})$ гёделева, то $\sem{\alpha}_{\Gamma(\mathcal{L})}=1$ или $\sem{\beta}_{\Gamma(\mathcal{L})}=1$.
\item Пусть, например, $\sem{\alpha}_{\Gamma(\mathcal{L})}=1$. Применим гомоморфизм: $\mathcal{G}(\sem{\alpha}_{\Gamma(\mathcal{L})})=\sem{\alpha}_{\mathcal{L}}=1_{\mathcal{L}}$.
\item $\sem{\alpha}_{\mathcal{L}}=[\alpha]=1_{\mathcal{L}}$ означает $\vd\alpha$ (билет~\hyperref[T:23]{23}, полнота $\mathcal{L}$).\qedhere
\end{enumerate}
\end{proof}

\textbf{Смысл.} Это формальный аналог BHK: чтобы доказать дизъюнкцию, надо уметь доказать один из дизъюнктов. Для КИВ неверно: $\vd_{\text{к}}A\vee\neg A$, но ни $A$, ни $\neg A$ не доказуемы. Дизъюнктивность поясняет замену аксиомы (10) на (10и).

\textbf{Следствия.} Закон исключённого третьего $\alpha\vee\neg\alpha$ недоказуем в ИИВ для переменной $A$ (иначе $\vd A$ или $\vd\neg A$, а оба не выводимы). Эта теорема --- второй (после моделей) способ показать недоказуемость ЗИТ.

\begin{warn}
\begin{itemize}
\item Теорема про \emph{доказуемость} без гипотез; с гипотезами неверно: $A\vee B\vd A\vee B$.
\item Для исчисления предикатов аналог: если $\vd\exists x\,\varphi(x)$, то есть терм $\theta$ с $\vd\varphi(\theta)$ (свойство существования) --- вне программы.
\end{itemize}
\end{warn}

\bigskip
\begin{center}
\fbox{\parbox{0.9\linewidth}{\centering Конец. Все билеты: \hyperref[map]{карта}\ \ $\cdot$\ \ \hyperref[toc]{оглавление}}}
\end{center}
\end{document}
